\documentclass[11pt]{article}
\usepackage{amsmath,amsthm,amssymb,mathtools}
\usepackage{bm}
\usepackage{hyperref}
\usepackage{cleveref}
\usepackage{geometry}
\usepackage{enumerate}
\geometry{margin=1in}

%% Theorem environments
\newtheorem{theorem}{Theorem}[section]
\newtheorem{lemma}[theorem]{Lemma}
\newtheorem{proposition}[theorem]{Proposition}
\newtheorem{corollary}[theorem]{Corollary}
\newtheorem{definition}[theorem]{Definition}
\newtheorem{assumption}[theorem]{Assumption}
\newtheorem{conjecture}[theorem]{Conjecture}
\newtheorem*{remark}{Remark}
\newtheorem*{example}{Example}

\DeclareMathOperator*{\argmax}{arg\,max}

\title{Coupled Exploration Collapse in Cooperative Contextual Bandits\\
\large (Corrected Version)}
\author{}
\date{}

\begin{document}
\maketitle
\tableofcontents
\newpage

%%============================================================
\section{Setting and Notation}\label{sec:setting}
%%============================================================

\paragraph{Communication graph.}
Let $G=(V,E)$ be a connected undirected graph with $N=|V|$ agents.
We associate to $G$ a symmetric doubly stochastic gossip matrix
$W\in\mathbb{R}^{N\times N}$ with $W_{ij}=0$ whenever $i\neq j$
and $(i,j)\notin E$.  The eigenvalues of $W$ are ordered as
$1=\lambda_1(W)>\lambda_2(W)\geq\cdots\geq\lambda_N(W)\geq -1$,
and we write $\gamma(W)\coloneqq 1-\lambda_2(W)\in(0,1]$ for the
spectral gap.  We also write
$\lambda_*(W)\coloneqq\max(|\lambda_2(W)|,|\lambda_N(W)|)<1$
for the second-largest eigenvalue \emph{in modulus}, which controls
the geometric decay of $W^k$ towards the uniform average.
The $\epsilon$-mixing time satisfies
$T_{\mathrm{mix}}(W;\epsilon)\leq(1-\lambda_*(W))^{-1}\log(1/\epsilon)$.

\begin{assumption}[Laziness / PSD gossip matrix]\label{ass:lazy}
$W$ is positive semidefinite, equivalently $\lambda_N(W)\geq 0$ and
hence $\lambda_*(W)=\lambda_2(W)$.  Any reversible $W$ can be made to
satisfy this by the standard lazy substitution
$W\mapsto\tfrac12(I+W)$, which preserves the stationary distribution
and at most doubles the mixing time.
\end{assumption}

\begin{remark}[Non-PSD case]\label{rem:lazy}
Assumption~\ref{ass:lazy} is used only to invoke per-step spectral
estimates (Lemmas~\ref{lem:gossip}, \ref{lem:leakage},
\ref{lem:prop}).  If $W$ is \emph{not} PSD (e.g.\ a Metropolis chain
with a negative eigenvalue), every per-$k$ bound below remains valid
after replacing $\lambda_2(W)^k$ by $\lambda_*(W)^k$, because
the relevant quantities are $|[W^k]_{ij}-1/N|$ and
$1-[W^k]_{ii}$, both of which are controlled by $\lambda_*$ rather
than $\lambda_2$.  We state the lemmas with $\lambda_*$ for this
reason; under Assumption~\ref{ass:lazy} one may read
$\lambda_*=\lambda_2$ throughout.
\end{remark}

\paragraph{Contextual social bandit.}
At each round $t=1,\dots,T$, every agent $i\in V$ observes a context
$\bm{x}_t^{(i)}\in\mathcal{X}\subset\mathbb{R}^d$ drawn i.i.d.\ from
an agent-specific distribution $\mathcal{D}_i$ (marginal $\mu_i$ on
$\mathcal{X}$, independent of history), selects an action
$a_t^{(i)}\in\mathcal{A}$ via a history-measurable policy, and
observes scalar feedback $y_t^{(i)}\in\mathbb{R}$.  For feature map
$\phi:\mathcal{X}\times\mathcal{A}\to\mathbb{R}^d$ and unknown
parameter $\bm{\theta}^{\!*}\in\mathbb{R}^d$:
\begin{equation}
y_t^{(i)}=\langle\phi(\bm{x}_t^{(i)},a_t^{(i)}),\bm{\theta}^{\!*}\rangle
+\beta^{(i)}\!\bigl(\bm{x}_t^{(i)},a_t^{(i)}\bigr)+\eta_t^{(i)},
\label{eq:observation_model}
\end{equation}
where $\beta^{(i)}:\mathcal{X}\times\mathcal{A}\to\mathbb{R}$ is
deterministic, agent-specific, bounded ($|\beta^{(i)}|\leq B_\beta$),
and $\eta_t^{(i)}$ is zero-mean $\sigma$-sub-Gaussian noise,
\emph{independent across $(t,i)$} conditional on
$\mathcal{F}_{t-1}$.  We write
$\bm{z}_t^{(i)}\coloneqq\phi(\bm{x}_t^{(i)},a_t^{(i)})$ and assume
$\|\bm{z}_t^{(i)}\|\leq L$ and $\|\bm{\theta}^{\!*}\|\leq S$.
Write $\pi^{\!*}(\bm{x})=\argmax_a\langle\phi(\bm{x},a),\bm{\theta}^{\!*}\rangle$
for the ground-truth optimal policy and measure performance via the
\emph{latent network regret}
\begin{equation}
R_T^{\mathrm{net}}\coloneqq\sum_{t=1}^T\sum_{i=1}^N
\Bigl[\langle\phi(\bm{x}_t^{(i)},\pi^{\!*}(\bm{x}_t^{(i)})),\bm{\theta}^{\!*}\rangle
-\langle\phi(\bm{x}_t^{(i)},a_t^{(i)}),\bm{\theta}^{\!*}\rangle\Bigr].
\label{eq:latent_regret}
\end{equation}

%%============================================================
\section{Context-Overlap and the Self-Consistent Bias Fixed Point}
\label{sec:overlap}
%%============================================================

\subsection{Policy-Dependent Bias Vector}

Fix $\bm{\theta}\in\mathbb{R}^d$ and define the greedy policy
$\pi_{\bm{\theta}}(\bm{x})\coloneqq\argmax_{a\in\mathcal{A}}
\langle\phi(\bm{x},a),\bm{\theta}\rangle$
with ties broken by a fixed measurable rule.  Let
\begin{align}
\bm{u}_i(\bm{\theta})&\coloneqq
\mathbb{E}_{\bm{x}\sim\mu_i}\bigl[\beta^{(i)}(\bm{x},\pi_{\bm{\theta}}(\bm{x}))\,
\phi(\bm{x},\pi_{\bm{\theta}}(\bm{x}))\bigr],\quad
\bar{\bm{u}}(\bm{\theta})\coloneqq\tfrac{1}{N}\sum_{i=1}^N\bm{u}_i(\bm{\theta}),
\label{eq:u_policy}\\
\Sigma_z^{(i)}(\bm{\theta})&\coloneqq
\mathbb{E}_{\bm{x}\sim\mu_i}\bigl[\phi(\bm{x},\pi_{\bm{\theta}}(\bm{x}))
\phi(\bm{x},\pi_{\bm{\theta}}(\bm{x}))^\top\bigr],\quad
\bar{\Sigma}_z(\bm{\theta})\coloneqq\tfrac{1}{N}\sum_{i=1}^N\Sigma_z^{(i)}(\bm{\theta}).
\label{eq:Sigma_policy}
\end{align}

\begin{definition}[Self-consistency map]\label{def:F}
Define $F:\mathbb{R}^d\to\mathbb{R}^d$ by
\[
F(\bm{\theta})\coloneqq\bm{\theta}^{\!*}+\bar{\Sigma}_z(\bm{\theta})^{-1}\bar{\bm{u}}(\bm{\theta}),
\qquad\text{wherever }\bar{\Sigma}_z(\bm{\theta})\succ 0.
\]
A \emph{biased fixed point} is any $\bm{\theta}_\infty$ satisfying
$F(\bm{\theta}_\infty)=\bm{\theta}_\infty$; equivalently,
$\bm{\delta}\coloneqq\bm{\theta}_\infty-\bm{\theta}^{\!*}$ solves
$\bar{\Sigma}_z(\bm{\theta}^{\!*}+\bm{\delta})\,\bm{\delta}
=\bar{\bm{u}}(\bm{\theta}^{\!*}+\bm{\delta})$.
\end{definition}

\subsection{Existence of a Biased Fixed Point}

\begin{assumption}[Regularity]\label{ass:reg}
(a)~Each marginal $\mu_i$ has a density with respect to Lebesgue
measure on $\mathcal{X}$.
(b)~There exist constants $R>0$ and $\sigma_{\min}>0$ such that
$\bar{\Sigma}_z(\bm{\theta})\succeq\sigma_{\min}I$ for all
$\bm{\theta}\in B_R\coloneqq\{\bm{\theta}:\|\bm{\theta}-\bm{\theta}^{\!*}\|\leq R\}$.
(c)~$\sigma_{\min}^{-1}\sup_{\bm{\theta}\in B_R}\|\bar{\bm{u}}(\bm{\theta})\|\leq R$
(self-mapping); this holds whenever $B_\beta L/\sigma_{\min}\leq R$.
\end{assumption}

\begin{assumption}[Margin / low-noise regularity]\label{ass:margin}
There is a constant $C_{\mathrm{mar}}<\infty$ such that for every pair
of actions $a\neq a'$, every $\bm{\theta}\in B_R$, and every $\zeta\geq 0$,
the $\mu_i$-mass of contexts whose decision margin is at most $\zeta$
is controlled,
\[
\mu_i\bigl(\{\bm{x}:\,
0<|\langle\phi(\bm{x},a)-\phi(\bm{x},a'),\bm{\theta}\rangle|\leq\zeta\}\bigr)
\leq C_{\mathrm{mar}}\,\zeta
\qquad(\forall i).
\]
\end{assumption}

\begin{remark}[Consequence: Lipschitz population maps]\label{rem:lip}
Assumption~\ref{ass:margin} is a Tsybakov-type low-noise/margin
condition.  Although the greedy selector $\pi_{\bm{\theta}}$ is a
\emph{discontinuous} function of $\bm{\theta}$, the condition forces
the $\mu_i$-mass of contexts that \emph{change} their selected action
when $\bm{\theta}$ moves by $\Delta\bm{\theta}$ to be
$O(C_{\mathrm{mar}}\cdot 2L\|\Delta\bm{\theta}\|)$ (a margin of size
$\leq 2L\|\Delta\bm{\theta}\|$ is needed to flip an action).  Since
the integrands defining $\bm{u}_i,\Sigma_z^{(i)}$ in
\eqref{eq:u_policy}--\eqref{eq:Sigma_policy} are bounded by $B_\beta L$
and $L^2$ respectively, the maps
$\bm{\theta}\mapsto\bar{\bm{u}}(\bm{\theta})$ and
$\bm{\theta}\mapsto\bar{\Sigma}_z(\bm{\theta})$ are Lipschitz on $B_R$,
with constants
\[
L_u\leq 2C_{\mathrm{mar}}B_\beta L^2,\qquad
L_\Sigma\leq 2C_{\mathrm{mar}}L^3.
\]
We invoke this in Lemma~\ref{lem:bias} (Step~4) and in
Lemma~\ref{lem:contraction}.  Without
Assumption~\ref{ass:margin} these maps can be discontinuous and the
fixed-point/perturbation arguments break down.
\end{remark}

For each $\bm{\theta}\in B_R$, let $\Pi(\bm{\theta})$ denote the set of
measurable policies $\pi:\mathcal{X}\to\mathcal{A}$ such that
$\pi(\bm{x})\in\argmax_a\langle\phi(\bm{x},a),\bm{\theta}\rangle$
for every $\bm{x}$.  For $\pi\in\Pi(\bm{\theta})$, define
\begin{align*}
\bar{\bm{u}}_\pi(\bm{\theta})&\coloneqq
\tfrac{1}{N}\sum_i\mathbb{E}_{\bm{x}\sim\mu_i}
\bigl[\beta^{(i)}(\bm{x},\pi(\bm{x}))\phi(\bm{x},\pi(\bm{x}))\bigr],\\
\bar{\Sigma}_{z,\pi}(\bm{\theta})&\coloneqq
\tfrac{1}{N}\sum_i\mathbb{E}_{\bm{x}\sim\mu_i}
\bigl[\phi(\bm{x},\pi(\bm{x}))\phi(\bm{x},\pi(\bm{x}))^\top\bigr].
\end{align*}

\begin{definition}[Set-valued self-consistency map]\label{def:Phi}
Define $\Phi:B_R\rightrightarrows B_R$ by
\[
\Phi(\bm{\theta})\coloneqq
\overline{\mathrm{co}}\,\bigl\{
\bm{\theta}^{\!*}+\bar{\Sigma}_{z,\pi}(\bm{\theta})^{-1}
\bar{\bm{u}}_\pi(\bm{\theta}):\pi\in\Pi(\bm{\theta})\bigr\},
\]
where $\overline{\mathrm{co}}$ denotes the closed convex hull.
\end{definition}

\begin{lemma}[Existence of biased fixed point]\label{lem:fp}
Under Assumption~\ref{ass:reg}, the set-valued map $\Phi$ has a fixed
point: there exists $\bm{\theta}_\infty\in B_R$ with
$\bm{\theta}_\infty\in\Phi(\bm{\theta}_\infty)$.
Equivalently, there exist $\bm{\theta}_\infty\in B_R$ and a
\emph{convex combination} of policies in $\Pi(\bm{\theta}_\infty)$---
meaning coefficients $\alpha_k\geq 0$, $\sum_k\alpha_k=1$, and
selectors $\pi_k\in\Pi(\bm{\theta}_\infty)$---such that
\begin{equation}
\bm{\theta}_\infty=\bm{\theta}^{\!*}+
\sum_k\alpha_k\,\bar{\Sigma}_{z,\pi_k}(\bm{\theta}_\infty)^{-1}
\bar{\bm{u}}_{\pi_k}(\bm{\theta}_\infty).
\label{eq:fp_conv}
\end{equation}
We write $\bm{\delta}\coloneqq\bm{\theta}_\infty-\bm{\theta}^{\!*}$.
\end{lemma}

\begin{proof}
We verify Kakutani's hypotheses for $\Phi:B_R\rightrightarrows B_R$.

\emph{(i) Nonemptiness and convexity of values.}  By the
Kuratowski--Ryll-Nardzewski measurable selection theorem,
$\Pi(\bm{\theta})\neq\emptyset$.
Assumption~\ref{ass:reg}(b) guarantees
$\bar{\Sigma}_{z,\pi}(\bm{\theta})\succeq\sigma_{\min}I$
for any $\pi\in\Pi(\bm{\theta})$, so inversion is well-defined.
Hence $\Phi(\bm{\theta})$ is nonempty; closed convex hulls are
nonempty, closed, and convex by construction.

\emph{(ii) Self-mapping.}  For any $\pi\in\Pi(\bm{\theta})$ and
$\bm{\theta}\in B_R$,
$\|\bar{\Sigma}_{z,\pi}(\bm{\theta})^{-1}\bar{\bm{u}}_\pi(\bm{\theta})\|
\leq\sigma_{\min}^{-1}B_\beta L\leq R$
by Assumption~\ref{ass:reg}(c).
Convex combinations of points in $B_R$ lie in $B_R$, so $\Phi:B_R\to B_R$.

\emph{(iii) Closed graph (upper hemicontinuity).}  Suppose
$\bm{\theta}_n\to\bm{\theta}$ in $B_R$, $\bm{y}_n\in\Phi(\bm{\theta}_n)$,
and $\bm{y}_n\to\bm{y}$.  We argue closedness of the graph of $\Phi$
directly, avoiding any (meaningless) ``almost-every sequence''
statement: the limiting point $\bm{\theta}$ is \emph{given}, not free
to be perturbed.

First, the value correspondence
$\bm{\theta}\mapsto\Pi(\bm{\theta})=\{\pi:\pi(\bm{x})\in
\argmax_a\langle\phi(\bm{x},a),\bm{\theta}\rangle\ \mu_i\text{-a.e.}\}$
is the (pointwise) argmax of the jointly continuous map
$(\bm{x},a,\bm{\theta})\mapsto\langle\phi(\bm{x},a),\bm{\theta}\rangle$
over the compact action set $\mathcal{A}$.  By Berge's maximum theorem
(see Aliprantis--Border, \emph{Infinite Dimensional Analysis},
Thm.~17.31), the pointwise argmax correspondence
$\bm{x}\rightrightarrows\argmax_a\langle\phi(\bm{x},a),\bm{\theta}\rangle$
is nonempty, compact-valued and \emph{upper hemicontinuous} in
$\bm{\theta}$, and admits a measurable selection by the
Kuratowski--Ryll-Nardzewski theorem (Aliprantis--Border,
Thm.~18.13).  This is what guarantees $\Pi(\bm{\theta})\neq\emptyset$ in
step~(i).

Second, we establish continuity of the integral functionals at the
given limit point $\bm{\theta}$.  Fix the selectors $\pi_n$ realizing
$\bm{y}_n$.  Decompose
$\mathcal{X}=\mathcal{X}^{\mathrm{tie}}_{\bm{\theta}}\cup
(\mathcal{X}\setminus\mathcal{X}^{\mathrm{tie}}_{\bm{\theta}})$, where
$\mathcal{X}^{\mathrm{tie}}_{\bm{\theta}}=\{\bm{x}:
\arg\max_a\langle\phi(\bm{x},a),\bm{\theta}\rangle
\text{ is not a singleton}\}$.
By Assumption~\ref{ass:margin} (with $\zeta\downarrow 0$),
$\mu_i(\mathcal{X}^{\mathrm{tie}}_{\bm{\theta}})=0$ for the given
$\bm{\theta}$ \emph{and every} $i$ --- this is precisely the role of
the margin condition, which removes the need for any ``generic
$\bm{\theta}$'' caveat.  On the complement, the argmax is unique and,
by upper hemicontinuity, locally constant along
$\bm{\theta}_n\to\bm{\theta}$; hence
$\beta^{(i)}(\bm{x},\pi_n(\bm{x}))\phi(\bm{x},\pi_n(\bm{x}))\to
\beta^{(i)}(\bm{x},\pi(\bm{x}))\phi(\bm{x},\pi(\bm{x}))$ pointwise for
$\mu_i$-a.e.\ $\bm{x}$, with $\pi(\bm{x})\in\argmax_a\langle\phi(\bm{x},a),
\bm{\theta}\rangle$.  The integrands are bounded ($\leq B_\beta L$ and
$\leq L^2$), so dominated convergence gives
$\bar{\bm{u}}_{\pi_n}(\bm{\theta}_n)\to\bar{\bm{u}}_\pi(\bm{\theta})$
and
$\bar{\Sigma}_{z,\pi_n}(\bm{\theta}_n)\to\bar{\Sigma}_{z,\pi}(\bm{\theta})$
with $\pi\in\Pi(\bm{\theta})$.
Continuity of matrix inversion on the positive-definite cone
(ensured by the $\sigma_{\min}I$ lower bound) then gives
convergence of the corresponding fixed-point candidates.
Since $\Phi(\bm{\theta})$ is the closed convex hull of such limits,
$\bm{y}\in\Phi(\bm{\theta})$, establishing the closed graph.

By Kakutani's theorem applied to the nonempty compact convex set $B_R$,
there exists $\bm{\theta}_\infty\in\Phi(\bm{\theta}_\infty)$.
Unpacking Definition~\ref{def:Phi}, $\bm{\theta}_\infty$ equals a
\emph{finite convex combination} of elements of the form
$\bm{\theta}^{\!*}+\bar{\Sigma}_{z,\pi_k}(\bm{\theta}_\infty)^{-1}
\bar{\bm{u}}_{\pi_k}(\bm{\theta}_\infty)$
for $\pi_k\in\Pi(\bm{\theta}_\infty)$,
which is the representation~\eqref{eq:fp_conv}.

\emph{Remark on representation.}
We do \emph{not} claim that the convex combination in~\eqref{eq:fp_conv}
reduces to a single term, because the map
$\pi\mapsto\bar{\Sigma}_{z,\pi}^{-1}\bar{\bm{u}}_\pi$ is nonlinear
(the matrix inverse is involved).  The representation as a convex
combination of selectors in $\Pi(\bm{\theta}_\infty)$ is the correct
and fully general form.  When the argmax is $\mu_i$-a.e.\ unique at
$\bm{\theta}_\infty$, all selectors agree and the combination reduces
to a single term.
\end{proof}

\begin{remark}[Non-uniqueness]\label{rem:nonunique}
The fixed point need not be unique.  Our negative result
(Theorem~\ref{thm:cec}) requires only that at least one
$\bm{\delta}\neq\bm{0}$ exists.
\end{remark}

\begin{assumption}[Non-trivial emergent bias]\label{ass:bias}
There exists a biased fixed point $\bm{\theta}_\infty\in B_R$ with
emergent bias
$\bm{\delta}=\bm{\theta}_\infty-\bm{\theta}^{\!*}$ satisfying
\[
\|\bm{\delta}\|_{\bar{\Sigma}_z(\bm{\theta}_\infty)}
\geq\rho_0\cdot\frac{B_\beta}{\sqrt{\lambda_{\max}(\bar{\Sigma}_z(\bm{\theta}_\infty))}}
\]
for some $\rho_0\in(0,1]$.
\end{assumption}


%%============================================================
\section{Naive Cooperative LinUCB and the Centralized Reference}
\label{sec:naive}
%%============================================================

\subsection{The Algorithm}

We analyze gossip-based cooperative LinUCB.  Each agent $i$ maintains
sufficient statistics $(A_t^{(i)},\bm{b}_t^{(i)})$ and at each round
performs a local update then a gossip step:
\begin{align*}
\widetilde{A}_t^{(i)}&=A_{t-1}^{(i)}+\bm{z}_t^{(i)}\bm{z}_t^{(i)\top},&
\widetilde{\bm{b}}_t^{(i)}&=\bm{b}_{t-1}^{(i)}+y_t^{(i)}\bm{z}_t^{(i)},\\
A_t^{(i)}&=\sum_{j=1}^N W_{ij}\,\widetilde{A}_t^{(j)},&
\bm{b}_t^{(i)}&=\sum_{j=1}^N W_{ij}\,\widetilde{\bm{b}}_t^{(j)},
\end{align*}
with $A_0^{(i)}=\lambda I$, $\bm{b}_0^{(i)}=\bm{0}$.
The estimator is
$\widehat{\bm{\theta}}_t^{(i)}=(A_t^{(i)})^{-1}\bm{b}_t^{(i)}$
and the action is
\[
a_t^{(i)}=\argmax_a\Bigl[\langle\phi(\bm{x}_t^{(i)},a),
\widehat{\bm{\theta}}_t^{(i)}\rangle
+\alpha_t\|\phi(\bm{x}_t^{(i)},a)\|_{(A_t^{(i)})^{-1}}\Bigr]
\]
with $\alpha_t=\sigma\sqrt{d\log\!\bigl(1+tL^2/(N\lambda)\bigr)+2\log(1/\delta)}
+\sqrt{\lambda}\,S$.

\subsection{The Centralized Reference Estimator}

\begin{definition}[Centralized reference]\label{def:cen}
\[
A_t^{\mathrm{cen}}\coloneqq\lambda I+\frac{1}{N}\sum_{\tau=1}^t\sum_{j=1}^N
\bm{z}_\tau^{(j)}\bm{z}_\tau^{(j)\top},\quad
\bm{b}_t^{\mathrm{cen}}\coloneqq\frac{1}{N}\sum_{\tau=1}^t\sum_{j=1}^N
y_\tau^{(j)}\bm{z}_\tau^{(j)},\quad
\bm{\theta}_t^{\mathrm{cen}}\coloneqq(A_t^{\mathrm{cen}})^{-1}\bm{b}_t^{\mathrm{cen}}.
\]
\end{definition}

\begin{lemma}[Gossip error bound]\label{lem:gossip}
Under the doubly-stochastic gossip iteration, for any agent $i$ and any $t\geq 1$:
\begin{equation}
\|A_t^{(i)}-A_t^{\mathrm{cen}}\|_{\mathrm{op}}
\leq\frac{\sqrt{N-1}\,L^2}{1-\lambda_*(W)},\qquad
\|\bm{b}_t^{(i)}-\bm{b}_t^{\mathrm{cen}}\|
\leq\frac{\sqrt{N-1}\,B_y L}{1-\lambda_*(W)},
\label{eq:gossip_err}
\end{equation}
where $B_y=SL+B_\beta+C\sigma\sqrt{\log(TN/\delta)}$
with high probability (uniform in $t$).
Under Assumption~\ref{ass:lazy} ($\lambda_*=\lambda_2$) the right-hand
sides read $\sqrt{N-1}\,L^2/\gamma(W)$ and
$\sqrt{N-1}\,B_y L/\gamma(W)$; in particular both scale as
$\Theta(\sqrt{N}/\gamma(W))$.
\end{lemma}

\begin{remark}[The $\sqrt{N}$ factor]\label{rem:sqrtN}
The cooperative--centralized gap grows like $\sqrt{N}$, not like a
constant.  A $k$-step walk started at the point mass $\delta_i$ carries
$\ell_1$ distance to uniform of order
$\sqrt{1/\pi_{\min}}\,\lambda_*^k=\sqrt{N}\,\lambda_*^k$ for the uniform
stationary law $\pi_{\min}=1/N$; the earlier constant $2\lambda_2^k$
understated this by exactly $\sqrt{N}$.  This $\sqrt{N}$ propagates
into $t^\star(G,\gamma)$ (Theorem~\ref{thm:cec}) and into term~(iii) of
Theorem~\ref{thm:dcesa}.
\end{remark}

\begin{proof}
\textbf{Unrolling the gossip recursion.}
By induction on the update rule, one verifies that
\[
A_t^{(i)}=\lambda I+\sum_{\tau=1}^t\sum_{j=1}^N
[W^{t-\tau+1}]_{ij}\,\bm{z}_\tau^{(j)}\bm{z}_\tau^{(j)\top},
\]
so that
\[
A_t^{(i)}-A_t^{\mathrm{cen}}
=\sum_{\tau=1}^t\sum_{j=1}^N
\Bigl([W^{t-\tau+1}]_{ij}-\tfrac{1}{N}\Bigr)
\bm{z}_\tau^{(j)}\bm{z}_\tau^{(j)\top}.
\]

\textbf{Key spectral bound.}
From the spectral decomposition of $W$,
$[W^k]_{ij}-1/N=\sum_{l\geq 2}\lambda_l^k v_{l,i}v_{l,j}$.
The sharp $\ell_1$ bound for a reversible chain started at the point
mass $i$ towards its (here uniform) stationary law is
\begin{equation}
\sum_{j=1}^N\bigl|[W^k]_{ij}-\tfrac{1}{N}\bigr|
=2\,\mathrm{TV}\!\bigl(W_i^{(k)},\mathrm{Uniform}\bigr)
\leq\sqrt{\frac{1-\pi_i}{\pi_i}}\;\lambda_*(W)^k
=\sqrt{N-1}\;\lambda_*(W)^k,
\label{eq:tv_bound}
\end{equation}
using $\pi_i=1/N$ and the standard spectral bound
$\mathrm{TV}\leq\tfrac12\sqrt{(1-\pi_i)/\pi_i}\,\lambda_*^k$ (e.g.\
Levin--Peres, \emph{Markov Chains and Mixing Times}, Prop.~12.19).
The modulus eigenvalue $\lambda_*=\max(|\lambda_2|,|\lambda_N|)$ is
required because, for non-PSD $W$, a negative $\lambda_N$ can dominate
at even/odd $k$; under Assumption~\ref{ass:lazy}, $\lambda_*=\lambda_2$.
The naive estimate $2\lambda_2^k$ omits the
$\sqrt{(1-\pi_i)/\pi_i}=\sqrt{N-1}$ factor carried by the point-mass
start (Remark~\ref{rem:sqrtN}).

\textbf{Operator-norm bound via triangle inequality.}
For each fixed $\tau$ (with $k=t-\tau+1$), applying the triangle
inequality and~\eqref{eq:tv_bound}:
\begin{align*}
\Bigl\|\sum_{j}\bigl([W^k]_{ij}-\tfrac{1}{N}\bigr)
\bm{z}_\tau^{(j)}\bm{z}_\tau^{(j)\top}\Bigr\|_{\mathrm{op}}
&\leq\sum_{j}\bigl|[W^k]_{ij}-\tfrac{1}{N}\bigr|\cdot
\|\bm{z}_\tau^{(j)}\bm{z}_\tau^{(j)\top}\|_{\mathrm{op}}\\
&\leq L^2\sum_{j}\bigl|[W^k]_{ij}-\tfrac{1}{N}\bigr|
\leq\sqrt{N-1}\,\lambda_*(W)^k L^2.
\end{align*}
Summing over $\tau=1,\dots,t$:
\[
\|A_t^{(i)}-A_t^{\mathrm{cen}}\|_{\mathrm{op}}
\leq\sqrt{N-1}\,L^2\sum_{k=1}^t\lambda_*(W)^k
\leq\frac{\sqrt{N-1}\,L^2\lambda_*(W)}{1-\lambda_*(W)}
\leq\frac{\sqrt{N-1}\,L^2}{1-\lambda_*(W)}.
\]

\textbf{Bound on $\bm{b}_t^{(i)}-\bm{b}_t^{\mathrm{cen}}$.}
The identical argument applies with
$\|y_\tau^{(j)}\bm{z}_\tau^{(j)}\|\leq B_y L$.
The bound $|y_\tau^{(j)}|\leq B_y$ a.s.\ with
$B_y\coloneqq SL+B_\beta+C\sigma\sqrt{\log(TN/\delta)}$
follows from $|y|\leq\|\phi\|\|\bm{\theta}^{\!*}\|+|\beta|+|\eta|$
and sub-Gaussian concentration of $\eta_\tau^{(j)}$ followed by a
union bound over $tN$ rounds.  The resulting bound is
$\|\bm{b}_t^{(i)}-\bm{b}_t^{\mathrm{cen}}\|\leq
\sqrt{N-1}\,B_y L/(1-\lambda_*(W))$.
\end{proof}

\subsection{Self-Normalized Concentration for the Centralized Estimator}

\begin{lemma}[Martingale concentration, centralized]\label{lem:mart}
Let $\bm{m}_t^{\mathrm{cen}}\coloneqq
N^{-1}\sum_{\tau=1}^t\sum_{j=1}^N\eta_\tau^{(j)}\bm{z}_\tau^{(j)}$.
With probability at least $1-\delta$, for all $t\geq 1$
simultaneously,
\begin{equation}
\|\bm{m}_t^{\mathrm{cen}}\|_{(A_t^{\mathrm{cen}})^{-1}}
\leq\frac{\beta_t}{\sqrt{N}},
\qquad
\beta_t\coloneqq\sigma\sqrt{d\log\!\bigl(1+tL^2/(N\lambda)\bigr)+2\log(1/\delta)}.
\label{eq:mart}
\end{equation}
(Here $\beta_t$ denotes the sub-Gaussian concentration radius; this
displayed form matches the derivation below exactly, with the
$2\log(1/\delta)$ confidence term kept separate from the
$d\log(1+\cdot)$ volumetric term, rather than the looser
$d\log((1+\cdot)/\delta)$ written previously.)
\end{lemma}

\begin{proof}
Enumerate agent-round pairs lexicographically:
for $s=(\tau-1)N+j$ define
$\epsilon_s\coloneqq\eta_\tau^{(j)}$
and $\bm{\xi}_s\coloneqq\bm{z}_\tau^{(j)}/\sqrt{N}$.
Then $\bm{\xi}_s$ is $\mathcal{G}_{s-1}$-measurable,
$\epsilon_s$ is $\sigma$-sub-Gaussian conditional on $\mathcal{G}_{s-1}$,
and $\|\bm{\xi}_s\|\leq L/\sqrt{N}$.
Setting $V_s\coloneqq\lambda I+\sum_{s'=1}^s\bm{\xi}_{s'}\bm{\xi}_{s'}^\top$
and $\bm{S}_s\coloneqq\sum_{s'=1}^s\epsilon_{s'}\bm{\xi}_{s'}$,
at $s=tN$ we have $V_{tN}=A_t^{\mathrm{cen}}$ and
$\bm{S}_{tN}=\sqrt{N}\,\bm{m}_t^{\mathrm{cen}}$.

Applying Theorem~1 of Abbasi-Yadkori, P\'al, and Szepesv\'ari~(2011):
with probability $\geq 1-\delta$, for all $s\geq 1$,
\[
\|\bm{S}_s\|_{V_s^{-1}}^2
\leq 2\sigma^2\log\!\left(\frac{\det(V_s)^{1/2}}
{\det(\lambda I)^{1/2}\cdot\delta}\right).
\]

To bound the log-determinant term at $s=tN$, we use the standard
identity (Abbasi-Yadkori et al.\ 2011, Lemma~10):
\[
\log\frac{\det(V_{tN})}{\det(\lambda I)}
=\log\det\!\Bigl(I+\tfrac{1}{\lambda}(V_{tN}-\lambda I)\Bigr)
\leq d\log\!\Bigl(1+\tfrac{\mathrm{tr}(V_{tN}-\lambda I)}{d\lambda}\Bigr).
\]
Since $\mathrm{tr}(V_{tN}-\lambda I)=N^{-1}\sum_{\tau,j}\|\bm{z}_\tau^{(j)}\|^2\leq tL^2$,
\[
\log\frac{\det(V_{tN})^{1/2}}{\det(\lambda I)^{1/2}}
\leq\frac{d}{2}\log\!\Bigl(1+\frac{tL^2}{N\lambda}\Bigr).
\]
Thus $\|\bm{S}_{tN}\|_{V_{tN}^{-1}}^2\leq
\sigma^2\!\left[d\log\!\bigl(1+tL^2/(N\lambda)\bigr)+2\log(1/\delta)\right]
=\beta_t^2$,
giving $\|\sqrt{N}\,\bm{m}_t^{\mathrm{cen}}\|_{(A_t^{\mathrm{cen}})^{-1}}\leq\beta_t$
and hence~\eqref{eq:mart}.
The time-uniform version uses the standard peeling argument in
Abbasi-Yadkori et al.
\end{proof}

\begin{lemma}[Contraction regime: uniqueness and geometric
population convergence]\label{lem:contraction}
Suppose Assumptions~\ref{ass:reg} and~\ref{ass:margin} hold, the
argmax is $\mu_i$-a.e.\ unique on $B_R$ (so $F$ of
Definition~\ref{def:F} is single-valued and, by
Remark~\ref{rem:lip}, Lipschitz), and the contraction modulus
\begin{equation}
\kappa\;\coloneqq\;\sigma_{\min}^{-1}\bigl(L_u+L_\Sigma R\bigr)
\;=\;\sigma_{\min}^{-1}\bigl(2C_{\mathrm{mar}}B_\beta L^2
+2C_{\mathrm{mar}}L^3 R\bigr)\;<\;1
\label{eq:kappa}
\end{equation}
(a joint margin / low-noise / small-bias condition: small
$C_{\mathrm{mar}}$, $B_\beta$, or large $\sigma_{\min}$).  Then $F$ is a
$\kappa$-contraction on $B_R$, it has a \emph{unique} fixed point
$\bm{\theta}_\infty\in B_R$, and the population iteration
$\bm{\theta}_{n+1}=F(\bm{\theta}_n)$ converges geometrically:
$\|\bm{\theta}_n-\bm{\theta}_\infty\|\leq\kappa^n\|\bm{\theta}_0-\bm{\theta}_\infty\|$.
\end{lemma}

\begin{proof}
For $\bm{\theta},\bm{\theta}'\in B_R$, write
$F(\bm{\theta})-F(\bm{\theta}')
=\bar{\Sigma}_z(\bm{\theta})^{-1}\bar{\bm{u}}(\bm{\theta})
-\bar{\Sigma}_z(\bm{\theta}')^{-1}\bar{\bm{u}}(\bm{\theta}')$.
Using the resolvent identity, the Lipschitz bounds of
Remark~\ref{rem:lip}
($\|\bar{\bm{u}}(\bm{\theta})-\bar{\bm{u}}(\bm{\theta}')\|\leq
L_u\|\bm{\theta}-\bm{\theta}'\|$,
$\|\bar{\Sigma}_z(\bm{\theta})-\bar{\Sigma}_z(\bm{\theta}')\|_{\mathrm{op}}
\leq L_\Sigma\|\bm{\theta}-\bm{\theta}'\|$),
$\bar{\Sigma}_z\succeq\sigma_{\min}I$, and
$\|\bar{\bm{u}}\|\leq B_\beta L\leq\sigma_{\min}R$,
\[
\|F(\bm{\theta})-F(\bm{\theta}')\|
\leq\sigma_{\min}^{-1}L_u\|\bm{\theta}-\bm{\theta}'\|
+\sigma_{\min}^{-2}L_\Sigma\|\bm{\theta}-\bm{\theta}'\|\cdot\|\bar{\bm{u}}\|
\leq\sigma_{\min}^{-1}(L_u+L_\Sigma R)\|\bm{\theta}-\bm{\theta}'\|
=\kappa\|\bm{\theta}-\bm{\theta}'\|.
\]
Banach's fixed-point theorem on the complete set $B_R$ (self-mapping by
Assumption~\ref{ass:reg}(c)) gives the conclusions.
\end{proof}

\begin{assumption}[Trajectory tracking; \emph{conditional}]\label{ass:pol_conv}
There exist a time $t_0<\infty$ and a deterministic sequence
$\{\epsilon_\tau\}_{\tau>t_0}$ with $\epsilon_\tau\to 0$ such that, on a
high-probability event (probability $\geq 1-\delta$, uniform in $\tau$),
\[
\|\bm{\theta}_\tau^{\mathrm{cen}}-\bm{\theta}_\infty\|\leq\epsilon_\tau
\quad(\tau>t_0),
\qquad\text{with partial sums}\quad
E_t\coloneqq\sum_{\tau=t_0+1}^t\epsilon_\tau=O(\sqrt{t}).
\]
\end{assumption}

\begin{remark}[On Assumption~\ref{ass:pol_conv}: what is and is not
proved]\label{rem:pol_conv}
We have \emph{weakened} the earlier hypothesis in two ways.
(i)~We no longer require \emph{summability}
$\sum_\tau\epsilon_\tau<\infty$: the natural statistical rate is
$\epsilon_\tau=\Theta(\sqrt{d\log(\tau/\delta)/(N\tau)})$ (from
Proposition~\ref{prop:conc}), whose partial sums are
$\Theta(\sqrt{t\log t/N})=O(\sqrt{t})$ but \emph{not} finite; the
weaker $E_t=O(\sqrt t)$ is exactly what
Lemma~\ref{lem:bias} consumes, so the strong summable form was both
unattainable and unnecessary.
(ii)~Lemma~\ref{lem:contraction} supplies a genuine
\emph{structural} basis: when $\kappa<1$ the fixed point is unique and
the \emph{population} map contracts geometrically, so there is a
single, well-defined target $\bm{\theta}_\infty$ and no fixed-point
multiplicity ambiguity.

What remains \emph{unproved} is the lift from population contraction to
the \emph{stochastic} trajectory: showing that the recursive-least-squares
iterate $\bm{\theta}_\tau^{\mathrm{cen}}$ (a running average, \emph{not}
a literal iteration of $F$) tracks $\bm{\theta}_\infty$ with
$E_t=O(\sqrt t)$ is a stochastic-approximation statement we do not
establish here.  We therefore keep Assumption~\ref{ass:pol_conv} as an
explicit hypothesis and state Theorem~\ref{thm:cec}
\emph{conditionally} on it.  The unconditional collapse theorem
(deriving Assumption~\ref{ass:pol_conv} from $\kappa<1$) is left open
(Caveat~C1).  Throughout we work in the basin of the unique
contraction fixed point $\bm{\theta}_\infty$ of
Lemma~\ref{lem:contraction}.
\end{remark}

\begin{lemma}[Bias concentration, centralized]\label{lem:bias}
Let $\bm{s}_t^{\mathrm{cen}}\coloneqq
N^{-1}\sum_{\tau=1}^t\sum_{j=1}^N
\beta^{(j)}(\bm{x}_\tau^{(j)},a_\tau^{(j)})\bm{z}_\tau^{(j)}$.
Under Assumptions~\ref{ass:margin} and~\ref{ass:pol_conv}, with
probability at least $1-\delta$, for all $t\geq t_0$:
\[
\Bigl\|\bm{s}_t^{\mathrm{cen}}-(A_t^{\mathrm{cen}}-\lambda I)\bm{\delta}\Bigr\|
\leq 2B_\beta L\sqrt{t\log(1/\delta)}+t_0 B_\beta L
+(B_\beta L+L^2\|\bm{\delta}\|)\,E_t,
\]
where $E_t=\sum_{\tau=t_0+1}^t\epsilon_\tau=O(\sqrt t)$ is the tracking
partial sum of Assumption~\ref{ass:pol_conv}.  In particular the entire
right-hand side is $O(\sqrt{t})$.
\end{lemma}

\begin{proof}
We decompose $\bm{s}_t^{\mathrm{cen}}$ into three parts:
\begin{align*}
\bm{s}_t^{\mathrm{cen}}&=
\underbrace{\frac{1}{N}\sum_{\tau=t_0+1}^t\sum_{j=1}^N
\bm{u}_j(\bm{\theta}_\tau^{\mathrm{cen}})}_{\mathrm{(A)\ expected\ drift}}
+\underbrace{\frac{1}{N}\sum_{\tau=t_0+1}^t\sum_{j=1}^N
\bigl[\bm{g}_\tau^{(j)}-\bm{u}_j(\bm{\theta}_\tau^{\mathrm{cen}})\bigr]}_{\mathrm{(B)\ martingale\ deviation}}
+\underbrace{\frac{1}{N}\sum_{\tau=1}^{t_0}\sum_{j=1}^N
\bm{g}_\tau^{(j)}}_{\mathrm{(C)\ pre-}t_0\mathrm{\ residual}},
\end{align*}
where $\bm{g}_\tau^{(j)}=\beta^{(j)}(\bm{x}_\tau^{(j)},a_\tau^{(j)})
\phi(\bm{x}_\tau^{(j)},a_\tau^{(j)})$.

\textit{Step 2 (C).}  Since $\|\bm{g}_\tau^{(j)}\|\leq B_\beta L$,
$\|(\mathrm{C})\|\leq t_0 B_\beta L$.

\textit{Step 3 (B).}  The increments
$\bm{g}_\tau^{(j)}-\bm{u}_j(\bm{\theta}_\tau^{\mathrm{cen}})$
have conditional mean zero and are bounded by $2B_\beta L$.
By the vector-valued Azuma--Hoeffding inequality (Hayes~2005,
Theorem~1.8), with probability $\geq 1-\delta$:
\[
\|(\mathrm{B})\|\leq 2B_\beta L\sqrt{\frac{(t-t_0)\log(2d/\delta)}{N}}
\leq 2B_\beta L\sqrt{t\log(1/\delta)},
\]
absorbing $d$ and constants into $\log(1/\delta)$.

\textit{Step 4 (A).}  Using the fixed-point relation
$\bar{\Sigma}_z(\bm{\theta}_\infty)\bm{\delta}=\bar{\bm{u}}(\bm{\theta}_\infty)$
and the tracking bound
$\|\bm{\theta}_\tau^{\mathrm{cen}}-\bm{\theta}_\infty\|\leq\epsilon_\tau$.
Here we crucially invoke Assumption~\ref{ass:margin}: by
Remark~\ref{rem:lip} it makes
$\bm{\theta}\mapsto\bar{\bm{u}}(\bm{\theta})$ and
$\bm{\theta}\mapsto\bar{\Sigma}_z(\bm{\theta})$ Lipschitz \emph{across}
the discontinuous greedy argmax (with constants
$L_u\leq 2C_{\mathrm{mar}}B_\beta L^2$,
$L_\Sigma\leq 2C_{\mathrm{mar}}L^3$), so that
$\|\bm{u}_j(\bm{\theta}_\tau^{\mathrm{cen}})-\bm{u}_j(\bm{\theta}_\infty)\|
\leq L_u\epsilon_\tau$ and
$\|(\bar{\Sigma}_z(\bm{\theta}_\tau^{\mathrm{cen}})
-\bar{\Sigma}_z(\bm{\theta}_\infty))\bm{\delta}\|\leq L_\Sigma\|\bm{\delta}\|\epsilon_\tau$.
Summing the per-round Lipschitz errors,
\[
\Bigl\|(\mathrm{A})-\sum_{\tau=t_0+1}^t\bar{\Sigma}_z(\bm{\theta}_\infty)\bm{\delta}\Bigr\|
\leq (B_\beta L+L^2\|\bm{\delta}\|)\sum_{\tau=t_0+1}^t\epsilon_\tau
=(B_\beta L+L^2\|\bm{\delta}\|)\,E_t,
\]
where $E_t=O(\sqrt t)$ by Assumption~\ref{ass:pol_conv}.  (Without
Assumption~\ref{ass:margin} the maps need not be Lipschitz and this
step fails.)

\textit{Step 5.}  By the matrix Azuma inequality (Tropp~2012,
Theorem~7.1), with probability $\geq 1-\delta$,
$A_t^{\mathrm{cen}}-\lambda I=\sum_\tau\bar{\Sigma}_z(\bm{\theta}_\tau^{\mathrm{cen}})+M_t^\Sigma$
where $\|M_t^\Sigma\|_{\mathrm{op}}\leq L^2\sqrt{2t\log(2d/\delta)}$.

\textit{Step 6 (Combining).}  On the intersection event
(probability $\geq 1-2\delta$):
\[
\|\bm{s}_t^{\mathrm{cen}}-(A_t^{\mathrm{cen}}-\lambda I)\bm{\delta}\|
\leq 2B_\beta L\sqrt{t\log(1/\delta)}+t_0 B_\beta L
+(B_\beta L+L^2\|\bm{\delta}\|)E_t
+L^2\|\bm{\delta}\|\sqrt{2t\log(2d/\delta)}.
\]
Since $L^2\|\bm{\delta}\|\sqrt{2t\log(2d/\delta)}$,
$(B_\beta L+L^2\|\bm{\delta}\|)E_t$ (with $E_t=O(\sqrt t)$), and
$2B_\beta L\sqrt{t\log(1/\delta)}$ are all $O(\sqrt{t})$,
we combine them into a single $O((B_\beta L+L^2\|\bm{\delta}\|)\sqrt{t\log(1/\delta)})$ term.
Rescaling $\delta\leftarrow\delta/2$ yields the stated
bound.
\end{proof}

\subsection{Concentration of the Centralized Estimator}

\begin{proposition}[Concentration around the biased fixed point]
\label{prop:conc}
Under Assumptions~\ref{ass:reg}, \ref{ass:margin}, \ref{ass:bias},
and~\ref{ass:pol_conv}, with probability $\geq 1-\delta$,
for all $t\geq t_0$, in the $A_t^{\mathrm{cen}}$-weighted norm:
\begin{equation}
\|\bm{\theta}_t^{\mathrm{cen}}-(\bm{\theta}^{\!*}+\bm{\delta})\|_{A_t^{\mathrm{cen}}}
\leq\underbrace{\frac{\beta_t}{\sqrt{N}}}_{\text{noise}}
+\underbrace{O\!\left(B_\beta L\sqrt{\tfrac{\log(1/\delta)}{\sigma_{\min}}}\right)}_{\Theta(1)\ \text{persistent bias}}
+\underbrace{\frac{\lambda\|\bm{\theta}^{\!*}+\bm{\delta}\|}{\sqrt{t\sigma_{\min}}}}_{\text{regularization}}.
\label{eq:conc}
\end{equation}
The bias term in the $A_t^{\mathrm{cen}}$-norm is $\Theta(1)$: it
does \emph{not} decay in $t$.  This is the crux of the result --- the
estimator is pulled a constant $A$-distance towards the biased fixed
point and stays there.
Equivalently, using the correct spectral scaling
$A_t^{\mathrm{cen}}\asymp t\,\bar{\Sigma}_z(\bm{\theta}_\infty)$ for large
$t$ (note: $\Theta(t)$, \emph{not} $\Theta(Nt)$; see
Remark~\ref{rem:spectral}), in the
$\bar{\Sigma}_z(\bm{\theta}_\infty)$-weighted norm
\begin{equation}
\|\bm{\theta}_t^{\mathrm{cen}}-(\bm{\theta}^{\!*}+\bm{\delta})\|_{\bar{\Sigma}_z(\bm{\theta}_\infty)}
\leq O\!\left(\frac{\beta_t}{\sqrt{N t}}\right)
+O\!\left(B_\beta L\sqrt{\tfrac{\log(1/\delta)}{t\sigma_{\min}}}\right)
+\frac{\lambda\|\bm{\theta}^{\!*}+\bm{\delta}\|}{t\sqrt{\sigma_{\min}}},
\label{eq:conc_sigma}
\end{equation}
where the noise term is $O(\beta_t/\sqrt{Nt})=O(\beta_t/\sqrt{Nt\sigma_{\min}})$.
The cooperative speedup is the $1/\sqrt N$ in the noise
\emph{variance} (because $\bm{m}_t^{\mathrm{cen}}$ averages $N$ agents),
\emph{not} an inflation of the design matrix to $\Theta(Nt)$.
\end{proposition}

\begin{remark}[Correct spectral scale of $A_t^{\mathrm{cen}}$]\label{rem:spectral}
The centralized design matrix carries a $1/N$ prefactor
(Definition~\ref{def:cen}) that exactly cancels the $N$ agents: its
trace satisfies
$\mathrm{tr}(A_t^{\mathrm{cen}}-\lambda I)
=N^{-1}\sum_{\tau,j}\|\bm{z}_\tau^{(j)}\|^2\leq tL^2$, so its eigenvalues
are $\Theta(t)$, and under Assumption~\ref{ass:reg}(b)
$A_t^{\mathrm{cen}}\succeq(\lambda+t\sigma_{\min})I\succeq t\sigma_{\min}I$
once $\frac1N\sum_j\bm{z}_\tau^{(j)}\bm{z}_\tau^{(j)\top}$ has accumulated
(formally, $A_t^{\mathrm{cen}}-\lambda I=\sum_{\tau}\bar{\Sigma}_z(\bm{\theta}_\tau^{\mathrm{cen}})+M_t^\Sigma$
by Lemma~\ref{lem:bias} Step~5, so $\lambda_{\min}(A_t^{\mathrm{cen}})\geq
t\sigma_{\min}-L^2\sqrt{2t\log(2d/\delta)}\geq\tfrac12 t\sigma_{\min}$ for
$t$ large).  In particular $\lambda_{\min}(A_t^{\mathrm{cen}})$
\emph{decreases} in $N$ for fixed per-round signal, and the bound
$A_t^{\mathrm{cen}}\succeq Nt\sigma_{\min}I$ used in an earlier draft was
off by a factor $N$.  Every downstream use of a $\Theta(Nt)$ design
matrix is corrected to $\Theta(t)$ below.
\end{remark}

\begin{proof}
Write
$\bm{\theta}_t^{\mathrm{cen}}-(\bm{\theta}^{\!*}+\bm{\delta})
=(A_t^{\mathrm{cen}})^{-1}\bigl[\bm{s}_t^{\mathrm{cen}}
-(A_t^{\mathrm{cen}}-\lambda I)\bm{\delta}
+\bm{m}_t^{\mathrm{cen}}-\lambda(\bm{\theta}^{\!*}+\bm{\delta})\bigr]$.

\emph{$A$-norm bound.}  By the triangle inequality in the
$(A_t^{\mathrm{cen}})^{-1}$-norm,
\[
\|\bm{\theta}_t^{\mathrm{cen}}-(\bm{\theta}^{\!*}+\bm{\delta})\|_{A_t^{\mathrm{cen}}}
\leq\underbrace{\|\bm{m}_t^{\mathrm{cen}}\|_{(A_t^{\mathrm{cen}})^{-1}}}_{\leq\beta_t/\sqrt N}
+\underbrace{\|\bm{s}_t^{\mathrm{cen}}-(A_t^{\mathrm{cen}}-\lambda I)\bm{\delta}\|_{(A_t^{\mathrm{cen}})^{-1}}}_{\text{bias}}
+\underbrace{\lambda\|\bm{\theta}^{\!*}+\bm{\delta}\|_{(A_t^{\mathrm{cen}})^{-1}}}_{\text{reg.}}.
\]
The noise term is $\leq\beta_t/\sqrt N$ by Lemma~\ref{lem:mart}.
For the bias term, $\|\bm{v}\|_{(A_t^{\mathrm{cen}})^{-1}}\leq
\|\bm{v}\|/\sqrt{\lambda_{\min}(A_t^{\mathrm{cen}})}\leq\|\bm{v}\|/\sqrt{t\sigma_{\min}/2}$
by Remark~\ref{rem:spectral}; with the Euclidean bias bound
$O(B_\beta L\sqrt{t\log(1/\delta)})$ of Lemma~\ref{lem:bias} this gives
\[
\|\bm{s}_t^{\mathrm{cen}}-(A_t^{\mathrm{cen}}-\lambda I)\bm{\delta}\|_{(A_t^{\mathrm{cen}})^{-1}}
\leq\frac{O(B_\beta L\sqrt{t\log(1/\delta)})}{\sqrt{t\sigma_{\min}/2}}
=O\!\left(B_\beta L\sqrt{\tfrac{\log(1/\delta)}{\sigma_{\min}}}\right),
\]
a quantity \emph{independent of $t$} --- the $\sqrt t$ in the
numerator and the $\sqrt t$ from $\lambda_{\min}(A_t^{\mathrm{cen}})$
cancel.  This is the persistent $\Theta(1)$ bias.
Similarly the regularization term is
$\leq\lambda\|\bm{\theta}^{\!*}+\bm{\delta}\|/\sqrt{t\sigma_{\min}/2}
=O(\lambda\|\bm{\theta}^{\!*}+\bm{\delta}\|/\sqrt{t\sigma_{\min}})$.
This proves~\eqref{eq:conc}.

\emph{$\bar{\Sigma}_z$-norm bound.}  By Remark~\ref{rem:spectral},
$A_t^{\mathrm{cen}}\asymp t\,\bar{\Sigma}_z(\bm{\theta}_\infty)$
($\Theta(t)$, not $\Theta(Nt)$).  Hence for any vector $\bm{w}$,
$\|\bm{w}\|_{\bar{\Sigma}_z}^2=\bm{w}^\top\bar{\Sigma}_z\bm{w}
\asymp t^{-1}\bm{w}^\top A_t^{\mathrm{cen}}\bm{w}
=t^{-1}\|\bm{w}\|_{A_t^{\mathrm{cen}}}^2$,
so $\|\bm{w}\|_{\bar{\Sigma}_z}\asymp t^{-1/2}\|\bm{w}\|_{A_t^{\mathrm{cen}}}$.
Applying this with
$\bm{w}=\bm{\theta}_t^{\mathrm{cen}}-(\bm{\theta}^{\!*}+\bm{\delta})$
and dividing each term of~\eqref{eq:conc} by $\sqrt t$ gives
\eqref{eq:conc_sigma}.  In particular the noise contribution is
$t^{-1/2}\cdot\beta_t/\sqrt N=\beta_t/\sqrt{Nt}$ --- the cooperative
gain is the $1/\sqrt N$, and the $1/\sqrt t$ is the usual estimation
decay; there is no spurious $1/N$ (the earlier
$\beta_t/(N\sqrt t)$ was the factor-$N$ error of
Remark~\ref{rem:spectral}).
\end{proof}

\begin{remark}[Collapse time]\label{rem:tstar}
The right-hand side of~\eqref{eq:conc_sigma} falls below
$\|\bm{\delta}\|_{\bar{\Sigma}_z(\bm{\theta}_\infty)}/4$
(equivalently~\eqref{eq:conc} below
$\|\bm{\delta}\|_{A_t^{\mathrm{cen}}}/4$, since the two norms differ by
the factor $\sqrt t$) for all $t\geq t_{\mathrm{cen}}^*$, where
\begin{equation}
t_{\mathrm{cen}}^*\coloneqq\Theta\!\left(
\underbrace{\frac{\sigma^2 d\log(T/\delta)}{N\,\|\bm{\delta}\|_{\bar{\Sigma}_z}^2}}_{\text{noise, }\propto 1/N}
+\underbrace{\frac{B_\beta^2 L^2\log(1/\delta)}{\sigma_{\min}\,\|\bm{\delta}\|_{\bar{\Sigma}_z}^2}}_{\text{bias, }\propto N^0}
\right).
\label{eq:tcen_def}
\end{equation}
The $1/N$ on the \emph{noise} term is the genuine cooperative speedup:
pooling $N$ agents reduces the noise term in~\eqref{eq:conc_sigma} from
$\beta_t/\sqrt t$ to $\beta_t/\sqrt{Nt}$, so the time at which it drops
below $\|\bm{\delta}\|_{\bar{\Sigma}_z}$ shrinks by $N$ (a $1/N$ law,
not the $1/N^2$ that the erroneous $\beta_t/(N\sqrt t)$ scaling would
have produced).  The \emph{bias} term carries no $1/N$: it reflects the
discrepancy between the running policy and the fixed point and is not
reduced by pooling.  Thus $t_{\mathrm{cen}}^*$ inherits the
$1/N$ scaling only when the noise term dominates.
\end{remark}


%%============================================================
\section{Coupled Exploration Collapse}\label{sec:cec}
%%============================================================

\begin{definition}[$\gamma$-decoupling region]\label{def:dec}
Fix an emergent bias $\bm{\delta}$ from Lemma~\ref{lem:fp}
and $\gamma>0$.
Let $\tilde{a}_{\bm{\delta}}(\bm{x})\coloneqq
\argmax_a\langle\phi(\bm{x},a),\bm{\theta}^{\!*}+\bm{\delta}\rangle$.
Define
\[
\mathcal{X}_{\mathrm{dec}}^\gamma\coloneqq\Bigl\{\bm{x}\in\mathcal{X}:\;
\langle\phi(\bm{x},\tilde{a}_{\bm{\delta}}(\bm{x}))
-\phi(\bm{x},\pi^{\!*}(\bm{x})),\;\bm{\theta}^{\!*}+\bm{\delta}\rangle\geq\gamma
\Bigr\}.
\]
The \emph{latent suboptimality} on this region is
\[
\Delta_{\min}^\gamma\coloneqq
\inf_{\bm{x}\in\mathcal{X}_{\mathrm{dec}}^\gamma}
\bigl[\langle\phi(\bm{x},\pi^{\!*}(\bm{x})),\bm{\theta}^{\!*}\rangle
-\langle\phi(\bm{x},\tilde{a}_{\bm{\delta}}(\bm{x})),\bm{\theta}^{\!*}\rangle\bigr]>0.
\]
Write $\mu_G=N^{-1}\sum_i\mu_i$.
\end{definition}

\begin{theorem}[Coupled Exploration Collapse; conditional on
Assumption~\ref{ass:pol_conv}]\label{thm:cec}
Fix $\gamma>0$ with $\mu_G(\mathcal{X}_{\mathrm{dec}}^\gamma)>0$,
$\Delta_{\min}^\gamma>0$, and the \emph{margin hypothesis}
\begin{equation}
\gamma\;\geq\;\frac{6\sqrt{2}\,\alpha_{t^\star}L}{\sqrt{t^\star\sigma_{\min}}}
\qquad(\text{note: }\sqrt{t^\star\sigma_{\min}},\ \emph{not}\ \sqrt{Nt^\star\sigma_{\min}}),
\label{eq:margin_hyp}
\end{equation}
where
\begin{equation}
t^\star(G,\gamma)\coloneqq
\underbrace{\Theta\!\left(
\frac{\sigma^2 dL^2\log(T/\delta)}{\gamma^2\sigma_{\min}}
+\frac{B_\beta^2 L^4\log(1/\delta)}{\gamma^2\sigma_{\min}^2}\right)}_{\text{estimate + per-agent bonus, no }1/N}
+\underbrace{\frac{C\sqrt{N}\,L^3(B_y+S)}{\gamma(W)\,\gamma\,\sigma_{\min}^{3/2}}}_{\text{gossip proximity, }\propto\sqrt N}.
\label{eq:tstar_full}
\end{equation}
Under Assumptions~\ref{ass:reg}, \ref{ass:margin}, \ref{ass:bias},
and (conditionally) \ref{ass:pol_conv}, naive cooperative LinUCB
satisfies, with probability $\geq 1-\delta$ (for
$N(T-t^\star)\mu_G(\mathcal{X}_{\mathrm{dec}}^\gamma)\geq C\log(1/\delta)$):
\begin{equation}
R_T^{\mathrm{net}}\geq
\tfrac12\,N\cdot\Delta_{\min}^\gamma\cdot\mu_G(\mathcal{X}_{\mathrm{dec}}^\gamma)
\cdot\bigl(T-t^\star(G,\gamma)\bigr).
\label{eq:cec_bound}
\end{equation}
In particular, for $T\gg t^\star(G,\gamma)$,
$R_T^{\mathrm{net}}=\Omega(NT)$.
\end{theorem}

\begin{remark}[The collapse time no longer enjoys a $1/N$ speedup in
its leading term]\label{rem:noN}
Contrast~\eqref{eq:tstar_full} with Remark~\ref{rem:tstar}: the
\emph{centralized estimator} concentrates at $t_{\mathrm{cen}}^*\propto
1/N$, but the \emph{collapse} requires in addition that each agent's
own UCB bonus---built from $A_t^{(i)}=\Theta(t)$, \emph{not}
$\Theta(Nt)$ (Remark~\ref{rem:spectral})---fall below $\gamma/3$.
That bonus condition gives the leading term of~\eqref{eq:tstar_full},
which carries \emph{no} $1/N$ (it is a factor $\sqrt N$ harder than the
erroneous $\sqrt{Nt^\star\sigma_{\min}}$ form would suggest: the bound is
on $\gamma\sqrt{t^\star}$, and dropping the spurious $\sqrt N$ raises the
required $t^\star$ by a factor $N$).  Pooling still accelerates the
\emph{estimate}; it does not shrink the per-agent exploration radius.
Consequently the cooperative speedup of collapse is weaker than the
earlier draft claimed.
\end{remark}

\begin{proof}
\textit{Step 1: Concentration around the biased fixed point.}
By Proposition~\ref{prop:conc} (bound~\eqref{eq:conc_sigma}) and the
choice of $t_{\mathrm{cen}}^*(\gamma)$ in
Remark~\ref{rem:tstar}/\eqref{eq:tstar_full}, for all $t$ at least the
first (no-$1/N$) bracket of~\eqref{eq:tstar_full}, with probability
$\geq 1-\delta$:
\[
\|\bm{\theta}_t^{\mathrm{cen}}-(\bm{\theta}^{\!*}+\bm{\delta})\|_{\bar{\Sigma}_z(\bm{\theta}_\infty)}
\leq\frac{\gamma\sqrt{\sigma_{\min}}}{12L}.
\]

By Lemma~\ref{lem:gossip}, using the resolvent identity
$(A_t^{(i)})^{-1}-(A_t^{\mathrm{cen}})^{-1}
=(A_t^{(i)})^{-1}(A_t^{\mathrm{cen}}-A_t^{(i)})(A_t^{\mathrm{cen}})^{-1}$
and the \emph{corrected} lower bound (Remark~\ref{rem:spectral})
$A_t^{(i)}\succeq A_t^{\mathrm{cen}}-\|A_t^{(i)}-A_t^{\mathrm{cen}}\|_{\mathrm{op}} I
\succeq(t\sigma_{\min}/2)I$
(for $t$ large enough that
$\sqrt{N}L^2/\gamma(W)\leq t\sigma_{\min}/2$, i.e.\
$t\geq 2\sqrt N L^2/(\gamma(W)\sigma_{\min})$), together with the
$\sqrt N$ gossip bound $\|A_t^{(i)}-A_t^{\mathrm{cen}}\|_{\mathrm{op}}\leq
\sqrt N L^2/\gamma(W)$ and the analogous $\bm{b}$ bound:
\[
\|\widehat{\bm{\theta}}_t^{(i)}-\bm{\theta}_t^{\mathrm{cen}}\|_{\bar{\Sigma}_z(\bm{\theta}_\infty)}
\leq\frac{C'\sqrt{N}\,L^2(B_y+S)}{\gamma(W)\,t\,\sigma_{\min}}.
\]
(The $\sqrt N$ comes from the gossip bound of Lemma~\ref{lem:gossip};
the single power of $t\sigma_{\min}$ replaces the earlier spurious
$Nt\sigma_{\min}^{1/2}$.)  For
$t\geq 12C'\sqrt N L^3(B_y+S)/(\gamma(W)\gamma\sigma_{\min}^{3/2})$
--- the second bracket of~\eqref{eq:tstar_full} --- this is
$\leq\gamma\sqrt{\sigma_{\min}}/(12L)$.
The triangle inequality then gives, for all $t\geq t^\star(G,\gamma)$:
\begin{equation}
\|\widehat{\bm{\theta}}_t^{(i)}-(\bm{\theta}^{\!*}+\bm{\delta})\|_{\bar{\Sigma}_z(\bm{\theta}_\infty)}
\leq\frac{\gamma\sqrt{\sigma_{\min}}}{6L}.
\label{eq:theta_hat_close}
\end{equation}

\textit{Step 2: Deterministic selection of the biased-optimal action.}
Fix $\bm{x}\in\mathcal{X}_{\mathrm{dec}}^\gamma$.
Write $\Delta_\phi=\phi(\bm{x},\tilde{a})-\phi(\bm{x},a^\star)$,
$L_\Delta=\|\Delta_\phi\|\leq 2L$.

The LinUCB score difference is
$\mathrm{UCB}_t^{(i)}(\tilde{a})-\mathrm{UCB}_t^{(i)}(a^\star)
=(\mathrm{I})+(\mathrm{II})$
where $(\mathrm{I})=\langle\Delta_\phi,\widehat{\bm{\theta}}_t^{(i)}\rangle$
and $(\mathrm{II})=\alpha_t[\|\phi(\bm{x},\tilde{a})\|_{(A_t^{(i)})^{-1}}
-\|\phi(\bm{x},a^\star)\|_{(A_t^{(i)})^{-1}}]$.

\textbf{Term (I).}
Writing $\widehat{\bm{\theta}}_t^{(i)}=\bm{\theta}^{\!*}+\bm{\delta}+\Delta\bm{\theta}$
with $\|\Delta\bm{\theta}\|_{\bar{\Sigma}_z}\leq\gamma\sqrt{\sigma_{\min}}/(6L)$
from~\eqref{eq:theta_hat_close}, and applying Cauchy--Schwarz
in the $\bar{\Sigma}_z$-weighted inner product
($\bar{\Sigma}_z\succeq\sigma_{\min}I$ gives
$\|\Delta_\phi\|_{\bar{\Sigma}_z^{-1}}\leq L_\Delta/\sqrt{\sigma_{\min}}$):
\[
(\mathrm{I})\geq\gamma-\frac{L_\Delta}{\sqrt{\sigma_{\min}}}
\cdot\frac{\gamma\sqrt{\sigma_{\min}}}{6L}\geq\gamma-\frac{2L}{6L}\gamma
=\frac{2\gamma}{3}.
\]

\textbf{Term (II).}
Since $A_t^{(i)}\succeq(t\sigma_{\min}/2)I$ (Remark~\ref{rem:spectral};
\emph{not} $Nt\sigma_{\min}/2$),
\[
|(\mathrm{II})|\leq 2\alpha_t\cdot\frac{L}{\sqrt{t\sigma_{\min}/2}}
=\frac{2\sqrt{2}\,\alpha_t L}{\sqrt{t\sigma_{\min}}}.
\]
The margin hypothesis~\eqref{eq:margin_hyp},
$\gamma\geq 6\sqrt{2}\,\alpha_{t^\star}L/\sqrt{t^\star\sigma_{\min}}$,
ensures this is $\leq\gamma/3$ for $t\geq t^\star$ (using that
$\alpha_t L/\sqrt{t\sigma_{\min}}$ is decreasing in $t$ up to
logarithmic factors).
(Note: the absence of $\sqrt N$ inside the square root makes this
condition a factor $\sqrt N$ \emph{harder} than the erroneous
$\sqrt{Nt^\star\sigma_{\min}}$ form: the per-agent design matrix is
$\Theta(t)$, so the exploration bonus does not shrink with $N$.  The
constant $6\sqrt 2\approx 8.49$ arises from the factor-of-$\sqrt 2$ in
the bonus bound.)

\textbf{Net score.}
$(\mathrm{I})+(\mathrm{II})\geq 2\gamma/3-\gamma/3=\gamma/3>0$.
Hence, for $t\geq t^\star$ and $\bm{x}_t^{(i)}\in\mathcal{X}_{\mathrm{dec}}^\gamma$,
the LinUCB rule selects $\tilde{a}_{\bm{\delta}}(\bm{x}_t^{(i)})$
rather than $\pi^{\!*}(\bm{x}_t^{(i)})$, deterministically.

\textit{Step 3: Accumulating latent regret (high probability).}
Each selection with $\bm{x}_t^{(i)}\in\mathcal{X}_{\mathrm{dec}}^\gamma$
(for $t>t^\star$) incurs instantaneous latent regret
$\geq\Delta_{\min}^\gamma$.  Let
$Z_t^{(i)}=\mathbb{1}[\bm{x}_t^{(i)}\in\mathcal{X}_{\mathrm{dec}}^\gamma]$;
these are independent across $(t,i)$ (contexts are drawn i.i.d.\
across agents and rounds) with
$\mathbb{E}[Z_t^{(i)}]=\mu_i(\mathcal{X}_{\mathrm{dec}}^\gamma)$, and
$\sum_{t>t^\star,i}\mathbb{E}[Z_t^{(i)}]
=N(T-t^\star)\mu_G(\mathcal{X}_{\mathrm{dec}}^\gamma)\eqqcolon m$.
By the multiplicative Chernoff bound,
$\Pr\bigl[\sum_{t>t^\star,i}Z_t^{(i)}<\tfrac12 m\bigr]\leq e^{-m/8}\leq\delta$
under the stated condition $m\geq 8\log(1/\delta)$.  On this event,
\[
R_T^{\mathrm{net}}\geq\Delta_{\min}^\gamma\sum_{t>t^\star,i}Z_t^{(i)}
\geq\tfrac12\,N\,\Delta_{\min}^\gamma\,\mu_G(\mathcal{X}_{\mathrm{dec}}^\gamma)\,(T-t^\star),
\]
which is~\eqref{eq:cec_bound}.  (Combining with the
probability-$\geq 1-\delta$ concentration event of Steps~1--2 and a
union bound rescales $\delta$ by a constant.)
\end{proof}

\begin{remark}[Why ``exploration collapse'']\label{rem:mech}
The two ingredients scale differently in $N$.  The \emph{estimator
error} towards the biased fixed point is $O(\beta_t/\sqrt{Nt})$
(Proposition~\ref{prop:conc}): cooperation \emph{accelerates} it by
$1/\sqrt N$.  The per-agent \emph{UCB bonus}, however, is
$\Theta(\alpha_t/\sqrt{t\sigma_{\min}})$, with \emph{no} $N$ speedup,
because $A_t^{(i)}=\Theta(t)$ (Remark~\ref{rem:spectral}); pooling does
not narrow the exploration radius.  Collapse occurs once \emph{both}
fall below $\Theta(\gamma)$: the biased point is identified quickly,
but exploration is suppressed only after the (slower, $N$-independent)
bonus shrinks.  Gossip contributes an additional
$\sqrt N/\gamma(W)$ proximity cost (Lemma~\ref{lem:gossip}).  Sparser
graphs (smaller $\gamma(W)$) delay collapse, giving the non-monotonicity
in $\gamma(W)$ (Section~\ref{sec:phase}).
\end{remark}


%%============================================================
\section{The Phase Transition}\label{sec:phase}
%%============================================================

Theorem~\ref{thm:cec} applies whenever $t^\star(G,\gamma)\ll T$,
which holds provided gossip mixes before the per-agent bonus has shrunk
below the margin.  Requiring each term of~\eqref{eq:tstar_full} to be
$\leq T/2$ yields the supercritical conditions
\begin{equation}
\gamma(W)\gtrsim\frac{\sqrt{N}\,L^3(B_y+S)}{T\,\gamma\,\sigma_{\min}^{3/2}},\qquad
\gamma^2\gtrsim\frac{\sigma^2 dL^2\log(T/\delta)}{T\sigma_{\min}}.
\label{eq:phase}
\end{equation}
The first is a \emph{gossip-rate} condition (now a factor $\sqrt N$
\emph{stricter}, and scaling as $1/\gamma$ rather than $1/\gamma^2$,
reflecting the $\sqrt N$ gossip bound of Lemma~\ref{lem:gossip}); the
second is a \emph{margin} condition (now a factor $N$ \emph{harder},
having lost the spurious $1/N$: the per-agent bonus enjoys no
cooperative speedup, Remark~\ref{rem:noN}).  Together they define the
supercritical regime where Theorem~\ref{thm:cec} gives
$R_T^{\mathrm{net}}=\Omega(NT)$.

\subsection{Individual Biased Fixed Points}\label{sec:indiv}

All machinery from Section~\ref{sec:overlap} specializes to $N=1$.
For each agent $i$, define
\[
\Phi_i(\bm{\theta})\coloneqq
\overline{\mathrm{co}}\bigl\{
\bm{\theta}^{\!*}+(\Sigma_{z,\pi}^{(i)}(\bm{\theta}))^{-1}
\bm{u}_{i,\pi}(\bm{\theta}):\pi\in\Pi(\bm{\theta})\bigr\}
\]
with expectations under $\bm{x}\sim\mu_i$ alone.
Kakutani applied to $\Phi_i$ on $B_R^{(i)}$ gives:

\begin{lemma}[Individual biased fixed point]\label{lem:fp_ind}
Under the single-agent analog of Assumption~\ref{ass:reg} for each $i$,
there exists $\bm{\theta}_\infty^{(i)}\in B_R^{(i)}$ with
$\bm{\theta}_\infty^{(i)}\in\Phi_i(\bm{\theta}_\infty^{(i)})$.
Write $\bm{\delta}_i\coloneqq\bm{\theta}_\infty^{(i)}-\bm{\theta}^{\!*}$.
\end{lemma}

\begin{assumption}[Individual bias non-triviality]\label{ass:bias_ind}
For each $i$, there exists a fixed point $\bm{\delta}_i$ satisfying
\[
\|\bm{\delta}_i\|_{\Sigma_z^{(i)}(\bm{\theta}_\infty^{(i)})}
\geq\rho_0^{\mathrm{ind}}\cdot
\frac{B_\beta}{\sqrt{\lambda_{\max}(\Sigma_z^{(i)}(\bm{\theta}_\infty^{(i)}))}}
\]
for some $\rho_0^{\mathrm{ind}}\in(0,1]$ (common across agents by
relabelling).
\end{assumption}

Define the \emph{individual} decoupling region at agent $i$:
\[
\mathcal{X}_{\mathrm{dec},i}^\gamma\coloneqq
\bigl\{\bm{x}\in\mathcal{X}:\;
\langle\phi(\bm{x},\tilde{a}_i(\bm{x}))-\phi(\bm{x},\pi^{\!*}(\bm{x})),
\bm{\theta}^{\!*}+\bm{\delta}_i\rangle\geq\gamma\bigr\},
\]
where $\tilde{a}_i(\bm{x})=\argmax_a\langle\phi(\bm{x},a),\bm{\theta}^{\!*}+\bm{\delta}_i\rangle$,
with latent gap
$\Delta_{\min,i}^\gamma=\inf_{\bm{x}\in\mathcal{X}_{\mathrm{dec},i}^\gamma}
[\langle\phi(\bm{x},\pi^{\!*}),\bm{\theta}^{\!*}\rangle
-\langle\phi(\bm{x},\tilde{a}_i),\bm{\theta}^{\!*}\rangle]>0$.

\subsection{Subcritical Leakage Bound}

\begin{definition}[Isolated reference estimator]
\[
A_t^{(i),\mathrm{iso}}\coloneqq\lambda I+\sum_{\tau=1}^t
\bm{z}_\tau^{(i)}\bm{z}_\tau^{(i)\top},\quad
\bm{b}_t^{(i),\mathrm{iso}}\coloneqq\sum_{\tau=1}^t
y_\tau^{(i)}\bm{z}_\tau^{(i)},\quad
\bm{\theta}_t^{(i),\mathrm{iso}}\coloneqq(A_t^{(i),\mathrm{iso}})^{-1}
\bm{b}_t^{(i),\mathrm{iso}}.
\]
\end{definition}

\begin{lemma}[Subcritical leakage bound]\label{lem:leakage}
Write $\gamma_\dagger(W)\coloneqq 1-\lambda_N(W)=\|I-W\|_{\mathrm{op}}$
for the \emph{full} spectral spread (so $\gamma_\dagger\geq\gamma(W)$,
with equality of order under Assumption~\ref{ass:lazy} when the
non-principal spectrum is one-sided).  For any $t\leq T$ and agent $i$:
\begin{equation}
\|A_t^{(i)}-A_t^{(i),\mathrm{iso}}\|_{\mathrm{op}}
\leq 2L^2\sum_{k=1}^t(1-[W^k]_{ii})
\leq 2L^2\,\gamma_\dagger(W)\sum_{k=1}^t k
\leq L^2\,\gamma_\dagger(W)\,t^2.
\label{eq:leakage_exact}
\end{equation}
In particular:
\begin{enumerate}[(a)]
\item \textbf{Deep subcritical bound (constant leakage):}
If $\gamma_\dagger(W)\leq c_0/T^2$ for an absolute constant $c_0>0$
(a near-isolation condition, implied by --- and in general stronger
than --- $\gamma(W)\leq c_0/T^2$), then for all $t\leq T$,
\begin{equation}
\|A_t^{(i)}-A_t^{(i),\mathrm{iso}}\|_{\mathrm{op}}
\leq L^2\gamma_\dagger(W)T^2\leq L^2 c_0.
\label{eq:leakage_const}
\end{equation}
\item \textbf{General bound:}
$\|A_t^{(i)}-A_t^{(i),\mathrm{iso}}\|_{\mathrm{op}}\leq 2L^2 t$
(trivial, from $\sum_{k=1}^t(1-[W^k]_{ii})\leq t$).
\end{enumerate}
The analogous bound for $\bm{b}_t^{(i)}-\bm{b}_t^{(i),\mathrm{iso}}$
holds with $L^2$ replaced by $B_y L$.
\end{lemma}

\begin{proof}
Unrolling the gossip recursion for the isolated comparison:
\[
A_t^{(i)}-A_t^{(i),\mathrm{iso}}
=\underbrace{\sum_{\tau=1}^t([W^{t-\tau+1}]_{ii}-1)
\bm{z}_\tau^{(i)}\bm{z}_\tau^{(i)\top}}_{\mathrm{(D1)}}
+\underbrace{\sum_{\tau=1}^t\sum_{j\neq i}[W^{t-\tau+1}]_{ij}
\bm{z}_\tau^{(j)}\bm{z}_\tau^{(j)\top}}_{\mathrm{(D2)}}.
\]

\textit{Bounding (D1) and (D2).}
Since $W$ is doubly stochastic and $[W^k]_{ij}\geq 0$,
$\sum_{j\neq i}[W^k]_{ij}=1-[W^k]_{ii}$.
Using $\|\bm{z}\bm{z}^\top\|_{\mathrm{op}}\leq L^2$:
\[
\|\mathrm{(D1)}\|_{\mathrm{op}}\leq L^2\sum_{k=1}^t(1-[W^k]_{ii}),
\qquad
\|\mathrm{(D2)}\|_{\mathrm{op}}\leq L^2\sum_{k=1}^t\sum_{j\neq i}[W^k]_{ij}
=L^2\sum_{k=1}^t(1-[W^k]_{ii}).
\]

\textit{Spectral bound on $1-[W^k]_{ii}$ (valid for non-PSD $W$).}
From the eigendecomposition $[W^k]_{ii}=1/N+\sum_{l\geq 2}\lambda_l^k v_{l,i}^2$
and $\sum_{l\geq 2}v_{l,i}^2=1-1/N$,
\[
1-[W^k]_{ii}=\sum_{l\geq 2}(1-\lambda_l^k)v_{l,i}^2
\leq\sum_{l\geq 2}k(1-\lambda_l)\,v_{l,i}^2
\leq k\,\gamma_\dagger(W)\sum_{l\geq 2}v_{l,i}^2
\leq k\,\gamma_\dagger(W),
\]
using the elementary inequality $1-\lambda^k\leq k(1-\lambda)$ for
\emph{every} $\lambda\in[-1,1]$ (so no positivity of $\lambda_l$ is
needed; this is the per-$k$ step where a naive $1-\lambda_2^k$ bound
would \emph{fail} at odd $k$ for a negative $\lambda_N$, cf.\
Remark~\ref{rem:lazy}), and $1-\lambda_l\leq 1-\lambda_N=\gamma_\dagger(W)$.
We deliberately use the bottom gap $\gamma_\dagger=1-\lambda_N$ rather
than $\gamma=1-\lambda_2$: the eigenvector-weighted sum
$\sum_{l\geq2}(\cdot)v_{l,i}^2$ can be dominated by \emph{any} non-top
mode, so it is $1-\lambda_N$, not $1-\lambda_2$, that rigorously
controls self-return.

\textit{Summing.}
$\sum_{k=1}^t(1-[W^k]_{ii})\leq\gamma_\dagger(W)\sum_{k=1}^t k
\leq\gamma_\dagger(W)T^2/2$, which proves~\eqref{eq:leakage_exact}.

\textit{Proof of part (a), constant leakage bound.}
For $\gamma_\dagger(W)\leq c_0/T^2$,
\[
\|A_t^{(i)}-A_t^{(i),\mathrm{iso}}\|_{\mathrm{op}}
\leq 2L^2\cdot\frac{\gamma_\dagger(W)T^2}{2}=L^2\gamma_\dagger(W)T^2\leq L^2 c_0.
\]
Since $\gamma_\dagger(W)\geq\gamma(W)$, the condition
$\gamma_\dagger(W)\leq c_0/T^2$ is a (possibly strictly stronger)
near-isolation condition than the spectral-gap condition
$\gamma(W)\leq c_0/T^2$; the two coincide up to constants whenever the
non-principal spectrum of $W$ is clustered near $\lambda_2$ (e.g.\ under
Assumption~\ref{ass:lazy} for a graph that is uniformly weakly
connected).
\end{proof}

\begin{remark}[Correct deep-subcritical threshold]
The $1/T^2$ scaling of the threshold is correct for constant leakage:
the key estimate $\sum_{k=1}^t(1-[W^k]_{ii})\leq\gamma_\dagger(W)T^2/2$
is valid whenever $T^2\gamma_\dagger(W)\leq c_0$, with no stricter power
of $T$ needed.  The one substantive correction relative to the earlier
draft is that the controlling spectral quantity is the \emph{bottom}
gap $\gamma_\dagger(W)=1-\lambda_N(W)=\|I-W\|_{\mathrm{op}}$, not the
\emph{spectral} gap $\gamma(W)=1-\lambda_2(W)$.  Operationally one reads
``$\gamma(W)\leq c_0/T^2$'' as the near-isolation condition
``$\|I-W\|_{\mathrm{op}}\leq c_0/T^2$''; in the deeply subcritical
(nearly disconnected, $W\approx I$) regime that motivates this lemma,
the whole non-principal spectrum sits near $1$ and the two readings
agree.
\end{remark}

\subsection{Subcritical Coupled Exploration Collapse}

\begin{proposition}[Individual concentration]\label{prop:conc_ind}
Under the single-agent forms of Assumptions~\ref{ass:reg}
and~\ref{ass:margin}, and Assumption~\ref{ass:bias_ind},
with probability $\geq 1-\delta/N$, for each $i$ and all $t\geq t_0$,
in the $A_t^{(i),\mathrm{iso}}$-norm (note $A_t^{(i),\mathrm{iso}}\succeq t\sigma_{\min}I$,
a $\Theta(t)$ scale --- the single-agent case carries no $N$ factor):
\[
\|\bm{\theta}_t^{(i),\mathrm{iso}}-(\bm{\theta}^{\!*}+\bm{\delta}_i)\|_{A_t^{(i),\mathrm{iso}}}
\leq\beta_t^{(i)}
+O\!\left(B_\beta L\sqrt{\tfrac{\log(N/\delta)}{\sigma_{\min}}}\right)
+\frac{\lambda\|\bm{\theta}^{\!*}+\bm{\delta}_i\|}{\sqrt{t\sigma_{\min}}},
\]
where $\beta_t^{(i)}=\sigma\sqrt{d\log(1+tL^2/\lambda)+2\log(N/\delta)}$
(matching the derivation, with the $2\log(N/\delta)$ confidence term
separate from the volumetric term).  Equivalently, in the
$\Sigma_z^{(i)}(\bm{\theta}_\infty^{(i)})$-norm every term is divided by
$\sqrt t$ (the $A$-norm bias term is $\Theta(1)$ and persistent, exactly
as in Proposition~\ref{prop:conc}; it is the $\Sigma_z^{(i)}$-norm that
vanishes).
The $\Sigma_z^{(i)}$-norm right-hand side falls below
$\gamma\sqrt{\sigma_{\min}}/(8L)$ for
$t\geq t_{\mathrm{iso}}^*(\gamma)$, where
\begin{equation}
t_{\mathrm{iso}}^*(\gamma)\coloneqq
\Theta\!\left(\frac{\sigma^2 dL^2\log(TN/\delta)+B_\beta^2 L^4\log(N/\delta)}
{\gamma^2\sigma_{\min}^2}\right).
\label{eq:tiso_star}
\end{equation}
\end{proposition}

\begin{theorem}[Subcritical CEC, individual failure mode]
\label{thm:cec_sub}
Suppose $\gamma_\dagger(W)\leq c_0/T^2$ for an absolute constant $c_0$
(Lemma~\ref{lem:leakage}(a)).
Fix $\gamma>0$ with $\min_i\mu_i(\mathcal{X}_{\mathrm{dec},i}^\gamma)>0$,
$\min_i\Delta_{\min,i}^\gamma>0$, and
$\gamma\geq 8\alpha_{t_{\mathrm{iso}}^*}L/\sqrt{t_{\mathrm{iso}}^*\sigma_{\min}}$
(no $N$ inside the root: the isolated design matrix is $\Theta(t)$).
Under Assumptions~\ref{ass:reg}, \ref{ass:margin}, and~\ref{ass:bias_ind},
naive cooperative LinUCB satisfies, with probability $\geq 1-\delta$:
\begin{equation}
R_T^{\mathrm{net}}\geq\sum_{i=1}^N\Delta_{\min,i}^\gamma\,
\mu_i(\mathcal{X}_{\mathrm{dec},i}^\gamma)\,(T-t_{\mathrm{iso}}^*(\gamma)).
\label{eq:cec_sub}
\end{equation}
For $T\gg t_{\mathrm{iso}}^*(\gamma)$, $R_T^{\mathrm{net}}=\Omega(NT)$.
\end{theorem}

\begin{proof}
\textit{Leakage bound.}
Fix agent $i$.  By Lemma~\ref{lem:leakage}(a) with
$\gamma_\dagger(W)\leq c_0/T^2$,
for all $t\leq T$:
$\|A_t^{(i)}-A_t^{(i),\mathrm{iso}}\|_{\mathrm{op}}\leq L^2 c_0$.

\textit{Estimator proximity.}
The resolvent identity gives
$(A_t^{(i)})^{-1}-(A_t^{(i),\mathrm{iso}})^{-1}
=(A_t^{(i)})^{-1}(A_t^{(i),\mathrm{iso}}-A_t^{(i)})(A_t^{(i),\mathrm{iso}})^{-1}$.
Using $\|(A_t^{(i)})^{-1}\|_{\mathrm{op}},\|(A_t^{(i),\mathrm{iso}})^{-1}\|_{\mathrm{op}}
\leq 1/(t\sigma_{\min})$:
\[
\|(A_t^{(i)})^{-1}-(A_t^{(i),\mathrm{iso}})^{-1}\|_{\mathrm{op}}
\leq\frac{L^2 c_0}{(t\sigma_{\min})^2}.
\]
The analogous bound for $\bm{b}$ gives
$\|\bm{b}_t^{(i)}-\bm{b}_t^{(i),\mathrm{iso}}\|\leq B_y L\,c_0$.
Combining:
\[
\|\widehat{\bm{\theta}}_t^{(i)}-\bm{\theta}_t^{(i),\mathrm{iso}}\|_{\Sigma_z^{(i)}}
\leq\frac{CL^3(B_y+S)c_0}{t^2\sigma_{\min}^{3/2}},
\]
which for $t\geq t_{\mathrm{iso}}^*(\gamma)$ is $o(\gamma\sqrt{\sigma_{\min}}/L)$
provided $T$ is large enough that
$CL^3(B_y+S)c_0/(t_{\mathrm{iso}}^*)^2\sigma_{\min}^{3/2}
\leq\gamma\sqrt{\sigma_{\min}}/(8L)$.\footnote{This gives a lower
bound on $T$: $T\geq C'L^4(B_y+S)c_0\sqrt{T}/(\sigma_{\min}^2\gamma)$,
i.e., $\sqrt{T}\geq C'L^4(B_y+S)c_0/(\sigma_{\min}^2\gamma)$,
a mild condition satisfied for large $T$.}

\textit{Proximity to individual fixed point.}
By Proposition~\ref{prop:conc_ind} and the leakage bound,
for $t\geq t_{\mathrm{iso}}^*(\gamma)$:
\[
\|\widehat{\bm{\theta}}_t^{(i)}-(\bm{\theta}^{\!*}+\bm{\delta}_i)\|_{\Sigma_z^{(i)}(\bm{\theta}_\infty^{(i)})}
\leq\frac{\gamma\sqrt{\sigma_{\min}}}{4L}
\]
with probability $\geq 1-\delta/N$.

\textit{Deterministic action selection.}
Applying the Step~2 argument of Theorem~\ref{thm:cec}'s proof
with $\widehat{\bm{\theta}}_t^{(i)}$ close to
$\bm{\theta}^{\!*}+\bm{\delta}_i$ instead of $\bm{\theta}^{\!*}+\bm{\delta}$,
agent $i$ selects $\tilde{a}_i(\bm{x}_t^{(i)})\neq\pi^{\!*}(\bm{x}_t^{(i)})$
on $\mathcal{X}_{\mathrm{dec},i}^\gamma$.

\textit{Regret accumulation.}
\[
R_T^{\mathrm{net}}
\geq\sum_{i=1}^N\sum_{t=t_{\mathrm{iso}}^*+1}^T
\mu_i(\mathcal{X}_{\mathrm{dec},i}^\gamma)\cdot\Delta_{\min,i}^\gamma.
\]
A union bound over $N$ agents gives overall probability $\geq 1-\delta$.
\end{proof}

\subsection{Unified Picture}

\begin{corollary}[Naive cooperative LinUCB fails in both regimes]
\label{cor:both}
Under Assumptions~\ref{ass:reg}, \ref{ass:bias}, \ref{ass:pol_conv},
and~\ref{ass:bias_ind}:
\begin{itemize}
\item \textbf{Supercritical,
$\gamma(W)\gtrsim\sqrt N L^3(B_y+S)/(T\gamma\sigma_{\min}^{3/2})$
and $\gamma^2\gtrsim\sigma^2 dL^2\log(T/\delta)/(T\sigma_{\min})$
(eq.~\eqref{eq:phase}):}
Theorem~\ref{thm:cec} (conditional on
Assumption~\ref{ass:pol_conv}) gives
$R_T^{\mathrm{net}}\geq \tfrac12 N\mu_G(\mathcal{X}_{\mathrm{dec}}^\gamma)
\Delta_{\min}^\gamma(T-t^\star(G,\gamma))$.
\item \textbf{Subcritical, $\gamma_\dagger(W)\leq c_0/T^2$:}
Theorem~\ref{thm:cec_sub} gives
$R_T^{\mathrm{net}}\geq\sum_i\mu_i(\mathcal{X}_{\mathrm{dec},i}^\gamma)
\Delta_{\min,i}^\gamma(T-t_{\mathrm{iso}}^*(\gamma))$
(unconditional --- it uses the isolated-reference analysis, not
Assumption~\ref{ass:pol_conv}).
\end{itemize}
Both are $\Omega(NT)$ for generic bias structures.  The two regimes
produce different biased policies: the supercritical bias is the
network-aggregate $\bm{\delta}$, while the subcritical bias is the
agent-local $\bm{\delta}_i$.
\end{corollary}

\begin{remark}[Intermediate regime]
Theorems~\ref{thm:cec} and~\ref{thm:cec_sub} cover the supercritical
regime~\eqref{eq:phase} and the deep-subcritical regime
$\gamma_\dagger(W)\leq c_0/T^2$ respectively.  The intermediate band is
not covered by either; we conjecture continuous interpolation between
the two bounds.
\end{remark}


%%============================================================
\section{Decentralized CESA}\label{sec:dcesa}
%%============================================================

\subsection{Algorithm}\label{sec:dcesa_algo}

Each agent $i$ maintains, for every $j\in\mathcal{N}(i)\cup\{i\}$,
a trust function
$w_{ij}(\bm{x},a)=\sigma(\langle\bm{\psi}_{ij},\phi(\bm{x},a)\rangle)$
with $\bm{\psi}_{ij}\in\mathbb{R}^d$ learned online.
Agent $i$ has private axiom access with per-round probability $p_i$.
When an axiom event occurs at round $t$ for agent $j$, it reveals
$z_t^{(j)}=\mathbb{1}[\langle\phi(\bm{x}_t^{(j)},a_t^{(j)}),
\bm{\theta}^{\!*}\rangle\geq\epsilon_{\mathrm{tol}}]$.
The trust-weighted cooperative update is
\begin{align*}
\widetilde{A}_t^{(i)}&=A_{t-1}^{(i)}+\bm{z}_t^{(i)}\bm{z}_t^{(i)\top},
&\widetilde{\bm{b}}_t^{(i)}&=\bm{b}_{t-1}^{(i)}+y_t^{(i)}\bm{z}_t^{(i)},\\
A_t^{(i)}&=\sum_j W_{ij}\,w_{ij}(\bm{x}_t^{(j)},a_t^{(j)})\,\widetilde{A}_t^{(j)},
&\bm{b}_t^{(i)}&=\sum_j W_{ij}\,w_{ij}(\bm{x}_t^{(j)},a_t^{(j)})\,\widetilde{\bm{b}}_t^{(j)}.
\end{align*}

\subsection{Axiom Propagation}\label{sec:axiom_prop}

At each round $t$, if agent $j$'s axiom fires, agent $j$ broadcasts
$\xi_t^{(j)}=(\phi(\bm{x}_t^{(j)},a_t^{(j)}),z_t^{(j)},y_t^{(j)})$
into a gossip buffer of depth $K\geq 1$.
After $K$ gossip rounds, agent $i$ has received $\xi_t^{(j)}$ with
probability proportional to $[W^K]_{ij}$.
Agent $i$ then performs one online-gradient step on $\bm{\psi}_{ij}$:
\[
\bm{\psi}_{ij}\leftarrow\bm{\psi}_{ij}
-\eta_t\nabla_{\bm{\psi}_{ij}}L_{\mathrm{CE}}\!\left(
\sigma(\langle\bm{\psi}_{ij},\phi(\bm{x}_t^{(j)},a_t^{(j)})\rangle),
\ell_t^{(i,j)}\right)
\]
with step-size $\eta_t=1/\sqrt{t}$ and $K=\lceil c/\gamma(W)\rceil$.

The \emph{effective axiom rate} at agent $i$ is
$p_t^{(i)}(K)\coloneqq\sum_{j\in V}[W^K]_{ij}p_j$,
with network average $\bar{p}=N^{-1}\sum_i p_i$.

\begin{lemma}[Effective axiom rate is network-uniform at mixing time]
\label{lem:prop}
For $K=\lceil c\log N/(1-\lambda_*(W))\rceil$ with $c\geq 1$
(under Assumption~\ref{ass:lazy}, $1-\lambda_*=\gamma(W)$), for all $i\in V$:
\[
\bigl|p_t^{(i)}(K)-\bar{p}\bigr|
\leq(N-1)\lambda_*(W)^K\leq N^{1-c}.
\]
For $c=2$, $|p_t^{(i)}(K)-\bar{p}|\leq 1/N$.
\end{lemma}

\begin{proof}
By the spectral decomposition and Cauchy--Schwarz
($\sum_{l\geq2}v_{l,i}^2\leq1$),
$|[W^K]_{ij}-1/N|=|\sum_{l\geq2}\lambda_l^K v_{l,i}v_{l,j}|\leq\lambda_*(W)^K$,
where the \emph{modulus} eigenvalue $\lambda_*=\max(|\lambda_2|,|\lambda_N|)$
is needed since a negative $\lambda_N$ would otherwise break the per-$K$
bound (Remark~\ref{rem:lazy}).  Hence
\[
\bigl|p_t^{(i)}(K)-\bar{p}\bigr|
=\Bigl|\sum_j([W^K]_{ij}-1/N)p_j\Bigr|
\leq(N-1)\lambda_*(W)^K.
\]
With $K=c\log N/(1-\lambda_*(W))$:
$\lambda_*(W)^K\leq e^{-K(1-\lambda_*(W))}=e^{-c\log N}=N^{-c}$.
\end{proof}

\begin{lemma}[Trust-classifier mistake bound]\label{lem:trust_mistakes}
Let $p^{(ij)}(K)=[W^K]_{ij}p_j$.  The expected cumulative
margin-mistakes $M_T^{(ij)}$ satisfies:
\[
\mathbb{E}[M_T^{(ij)}]
\leq\frac{C_1\sqrt{Td\log T}}{\Delta_0\sqrt{p^{(ij)}(K)}}
+\frac{C_2\log(1/\delta)}{\Delta_0^2\,p^{(ij)}(K)}.
\]
\end{lemma}

\begin{proof}
Three parts: (1)~OGD regret on labeled stream; (2)~loss-to-margin
conversion; (3)~extrapolation to unlabeled rounds via Bernstein.

\textit{Part 1.}  With $S_T^{(ij)}$ labeled rounds, the OGD
regret bound (Hazan~2016, Theorem~3.1) gives
$\sum_{s=1}^{S_T^{(ij)}}\ell(\bm{\psi}_{ij}^{(s)};\cdot)
\leq D_\psi L\sqrt{S_T^{(ij)}}$
under the realizability assumption.

\textit{Part 2.}  A margin-mistake $E_s$ implies
$\ell(\bm{\psi}_{ij}^{(s)};\cdot)\geq\Delta_0$, so the count of
labeled margin-mistakes is $\leq D_\psi L\sqrt{S_T^{(ij)}}/\Delta_0$.

\textit{Part 3.}  The unlabeled-mistake expected count is bounded by
$D_\psi L\sqrt{T}/(\Delta_0\sqrt{p^{(ij)}(K)})$ via online-to-batch
conversion.  Bernstein's inequality (Boucheron--Lugosi--Massart~2013,
Theorem~2.10) then gives the stated high-probability bound.
The $\sqrt{d\log T}$ factor arises from a covering-number
argument over the classifier space.
\end{proof}

\subsection{D-CESA Regret Bound}\label{sec:dcesa_regret}

\begin{theorem}[D-CESA network regret with axiom propagation]
\label{thm:dcesa}
Under Assumption~\ref{ass:reg}, assuming $\mathcal{H}_{\mathrm{trust}}$
realizes the $\epsilon_{\mathrm{tol}}$-indicator for each neighbor,
and using $K=\lceil c\log N/\gamma(W)\rceil$ with $c\geq 2$:
\begin{equation}
R_T^{\mathrm{net}}\leq
\underbrace{\widetilde{O}\!\left(|E|\sqrt{\frac{Td}{\bar{p}}}\right)}_{(i)\text{ epistemic}}
+\underbrace{\widetilde{O}\!\left(d\sqrt{NT}\right)}_{(ii)\text{ cooperative bandit}}
+\underbrace{\widetilde{O}\!\left(\frac{d\sqrt{NT}}{\gamma(W)}\right)}_{(iii)\text{ gossip error}}
+\underbrace{\epsilon_{\mathrm{tol}}\,NT}_{(iv)\text{ tolerance}}.
\label{eq:dcesa}
\end{equation}
Term~(i) depends on the network-average rate $\bar{p}$, not $\min_i p_i$.
Term~(iii) carries a $\sqrt N$ relative to the earlier draft, inherited
from the corrected $\sqrt N$ gossip bound of Lemma~\ref{lem:gossip}
(Remark~\ref{rem:sqrtN}).
\end{theorem}

\begin{proof}
Decompose $R_T^{\mathrm{net}}$ into trust-mistake rounds (I),
gossip-error rounds (III), and clean rounds (II+IV).

\textit{Term (I).}
On trust-mistake rounds, per-round regret is at most 1.
Lemma~\ref{lem:trust_mistakes} summed over edges, combined with
Lemma~\ref{lem:prop} (giving $p^{(ij)}(K)\geq\bar{p}/2$ for $N\geq 2$),
yields $(\mathrm{I})\leq\widetilde{O}(|E|\sqrt{Td/\bar{p}})$.

\textit{Term (II+IV).}
On clean rounds, the trust-weighted centralized estimator satisfies
a weighted analog of Lemma~\ref{lem:mart} (with
$\bm{\xi}_s=\sqrt{w_{ij}/N}\,\bm{z}_s$, preserving the
Abbasi-Yadkori martingale hypotheses since weights $w_{ij}$ are
$\mathcal{G}_{s-1}$-measurable and $\epsilon_s$ remains
$\sigma$-sub-Gaussian).  The resulting concentration bound gives
$\alpha_t^w\leq\alpha_t$.  Following the LinUCB regret analysis with
the elliptical potential lemma and Cauchy--Schwarz:
$(\mathrm{II+IV})\leq\widetilde{O}(d\sqrt{NT})+\epsilon_{\mathrm{tol}}NLT$.

\textit{Term (III).}
The gossip deviation
$\|A_t^{(i)}-A_t^{\mathrm{cen,w}}\|_{\mathrm{op}}\leq\sqrt{N-1}\,L^2/(1-\lambda_*(W))
=\widetilde{O}(\sqrt N L^2/\gamma(W))$
(Lemma~\ref{lem:gossip}, with the corrected $\sqrt N$ factor).
Propagating through the resolvent identity and summing over $T$ rounds
gives $(\mathrm{III})\leq\widetilde{O}(\sqrt N\,d\sqrt{T}/\gamma(W))
=\widetilde{O}(d\sqrt{NT}/\gamma(W))$.
\end{proof}


%%============================================================
\section{Lower Bounds for D-CESA}\label{sec:dcesa_lower}
%%============================================================

Lower bounds are stated against any cooperative bandit algorithm
with heterogeneous axiom access; they are information-theoretic.

\subsection{Lower Bound for the Epistemic Term (i)}

\begin{assumption}[Soft / noisy certification channel]\label{ass:softcert}
The certification observed on an axiom-fired round is a \emph{noisy}
label: instead of the deterministic indicator
$\mathbb{1}[\langle\phi,\bm{\theta}^{\!*}\rangle\geq\epsilon_{\mathrm{tol}}]$,
the learner observes $z=\langle\phi,\bm{\theta}^{\!*}\rangle+\zeta$ with
$\zeta$ independent, zero-mean, $\sigma_z$-sub-Gaussian (equivalently a
logistic/Bernoulli soft label with margin scale $\sigma_z$).  This is
consistent with the \emph{learned} soft trust classifier of
Section~\ref{sec:dcesa_algo}, whose output is a probability, not a hard
bit.
\end{assumption}

\begin{theorem}[Epistemic lower bound; conditional on
Assumption~\ref{ass:softcert}]\label{thm:lower_epistemic}
Under Assumption~\ref{ass:softcert} there exists a family of problem
instances on which any algorithm $\mathcal{A}$ with heterogeneous axiom
access $(p_i)_{i\in V}$ and effective network rate $\bar{p}$ incurs
\[
\mathbb{E}[R_T^{\mathrm{net}}]\geq\Omega\!\left(\sigma_z|E|\sqrt{Td/\bar{p}}\,\right).
\]
Without Assumption~\ref{ass:softcert} (i.e.\ for the noiseless
deterministic axiom of Section~\ref{sec:dcesa_algo}), only the weaker
waiting-time bound $\mathbb{E}[R_T^{\mathrm{net}}]\geq\Omega(|E|\,\Delta_{\mathrm{gap}}/\bar p)$
holds (Remark~\ref{rem:epi_downgrade}).
\end{theorem}

\begin{proof}
The construction is rebuilt to repair two defects of an earlier
version: (a)~it used $\bm{\theta}^{\!*}=0$, which makes
$R_T^{\mathrm{net}}\equiv 0$ and lower-bounds nothing; and (b)~it
claimed zero KL on non-axiom rounds even though the signed reward
$y=\langle\phi,\bm{\theta}^{\!*}\rangle+\beta+\eta$ has a
\emph{bias-dependent mean}, so raw feedback \emph{does} leak the
discriminating direction.  We fix both by making the bias exactly
cancel the reward gap, so reward is genuinely uninformative, while
$\bm{\theta}^{\!*}\neq 0$ produces real latent regret.

\textit{Step 1: Construction with $\bm{\theta}^{\!*}\neq 0$ and reward
cancellation.}
Fix a directed edge $(i,j)$ and a coordinate $e_k$.  Two actions
$a_+,a_-$ have fixed features $\phi(\bm{x},a_\pm)=\pm e_k$.  In instance
$(k,s)$, $s\in\{+,-\}$, set
$\bm{\theta}^{\!*}_{(k,s)}=s\,\Delta\,e_k$, so under $s=+$ action $a_+$
is truly optimal (latent value $+\Delta$ vs $-\Delta$) and under $s=-$
action $a_-$ is.  Choose the bias to cancel the latent reward:
$\beta^{(j)}_{(k,s)}(\bm{x},a)=-\langle\phi(\bm{x},a),\bm{\theta}^{\!*}_{(k,s)}\rangle$,
so the \emph{observed} reward mean
$\langle\phi(\bm{x},a),\bm{\theta}^{\!*}_{(k,s)}\rangle
+\beta^{(j)}_{(k,s)}(\bm{x},a)\equiv 0$
for every action and every $s$ (and $|\beta|\leq\Delta\leq B_\beta$).

\textit{Step 2: Reward is exactly uninformative.}
By Step~1 the reward is $y=0+\eta$, with law $N(0,\sigma^2)$ (or the
common sub-Gaussian law) \emph{identical} across actions and across
$s\in\{+,-\}$.  Hence the per-round KL contributed by the reward
channel is genuinely zero:
$\mathrm{KL}(\mathbb{P}_{(k,+)}^{\mathrm{reward}}\|
\mathbb{P}_{(k,-)}^{\mathrm{reward}})=0$ on every round, axiom or not.
This is now a \emph{correct} statement (it was false in the earlier
version, where the bias entered the reward mean).

\textit{Step 3: KL from the (noisy) axiom channel.}
By Assumption~\ref{ass:softcert}, on an axiom round at agent $j$ the
learner observes $z=\langle\phi(\bm{x},a),\bm{\theta}^{\!*}_{(k,s)}\rangle+\zeta
=s\,\Delta\,\langle\phi(\bm{x},a),e_k\rangle+\zeta$.  For action $a_+$
the mean of $z$ is $+s\Delta$, so the two hypotheses differ in the
axiom-channel mean by $2\Delta$, giving per-axiom-round KL
$\leq 2\Delta^2/\sigma_z^2$.  Effective axiom rate at agent $i$,
summed over sources $j$, is $\bar p$ per round
(Lemma~\ref{lem:prop}).  By the KL chain rule over the $\approx T\bar p$
axiom rounds,
\[
\mathrm{KL}(\mathbb{P}_{(k,+)}^T\|\mathbb{P}_{(k,-)}^T)
=\underbrace{0}_{\text{reward}}+\underbrace{C_{\mathrm{KL}}\,\frac{\Delta^2}{\sigma_z^2}\,T\bar p}_{\text{axiom}}.
\]
The factor $\bar p$ (not $N$) is now legitimate: the reward channel
contributes exactly $0$ (Step~2), so \emph{only} axiom rounds carry
information, and their total rate per round is $\bar p$.

\textit{Step 4: Le Cam.}  Distinguishing the two hypotheses with
probability $\geq 3/4$ requires (Pinsker) $\mathrm{KL}\geq c_0$, i.e.\
$\Delta\geq c\,\sigma_z/\sqrt{T\bar p}$.  Since reward gives no signal,
any algorithm that has not distinguished them pulls the wrong action a
constant fraction of the $T$ rounds, incurring expected latent regret
$\Omega(\Delta T)$ on at least one of the pair; optimizing
$\Delta\asymp\sigma_z/\sqrt{T\bar p}$ gives $\Omega(\sigma_z\sqrt{T/\bar p})$
per edge.

\textit{Step 5: Assouad over $d$ coordinates.}
Embedding $d$ orthogonal coordinate constructions and applying
Assouad's lemma yields $\Omega(\sigma_z\sqrt{Td/\bar p})$ per directed
edge.

\textit{Step 6: Summing over edges.}
The per-edge action subsets $\mathcal{A}_{(i,j)}$ are disjoint, so the
sub-problems are independent and the regrets add:
$\mathbb{E}[R_T^{\mathrm{net}}]\geq\Omega(\sigma_z|E|\sqrt{Td/\bar p})$.
\qedhere
\end{proof}

\begin{remark}[Why a noise assumption is needed, and the noiseless
downgrade]\label{rem:epi_downgrade}
The $\sqrt{\,\cdot\,}$ rate is an estimation rate and \emph{requires} a
noisy discriminating channel.  Under the literal noiseless axiom
$z=\mathbb{1}[\langle\phi,\bm{\theta}^{\!*}\rangle\geq\epsilon_{\mathrm{tol}}]$
of Section~\ref{sec:dcesa_algo}, a single axiom observation of $a_+$
reveals $s$ \emph{exactly}; identification then costs only the
geometric waiting time $\Theta(1/\bar p)$ for the first axiom, giving
the weaker bound $\Omega(|E|\,\Delta_{\mathrm{gap}}/\bar p)$ (no
$\sqrt{T}$, no $\sqrt d$ from repeated estimation).
Assumption~\ref{ass:softcert} is what converts the problem into a
genuine estimation problem with the $\sqrt{Td/\bar p}$ rate.  We do not
claim the unconditional $\sqrt{Td/\bar p}$ lower bound for the
deterministic-axiom model; closing this gap (matching the
upper bound term~(i) for the hard-label model) is left \emph{open}.
\end{remark}

\begin{remark}[Tightness: partial]\label{rem:epi_tight}
Term~(i) of Theorem~\ref{thm:dcesa} is
$\widetilde{O}(|E|\sqrt{Td/\bar{p}})$.  Under
Assumption~\ref{ass:softcert} (soft labels, the regime in which the
trust \emph{classifier} of Section~\ref{sec:dcesa} actually operates),
Theorem~\ref{thm:lower_epistemic} matches it up to the $\sigma_z$ scale
and logarithmic factors, so the $\bar p$-dependence is necessary
\emph{in that regime}.  For the hard-label model the matching lower
bound is open (Remark~\ref{rem:epi_downgrade}); the summary table marks
term~(i) tightness accordingly.
\end{remark}

\subsection{Lower Bound for the Cooperative Bandit Term (ii)}

\begin{theorem}[Cooperative bandit lower bound]\label{thm:lower_coop}
There exists a family of linear cooperative bandit instances
(no bias, no trust learning) on which any cooperative algorithm
$\mathcal{A}$ satisfies:
\[
\mathbb{E}[R_T^{\mathrm{net}}]\geq\Omega(d\sqrt{NT}).
\]
\end{theorem}

\begin{proof}
Set $\beta^{(i)}=0$ and $p_i=1$.  The problem reduces to cooperative
LinUCB with unbiased feedback.  The $\Omega(d\sqrt{NT})$ lower bound
for network regret follows from Dubey--Pentland (2020, Theorem~2)
adapted to linear bandits.
\end{proof}

\subsection{Lower Bound for the Tolerance Term (iv)}

\begin{theorem}[Tolerance lower bound]\label{thm:lower_tol}
If trust certifications are $\epsilon_{\mathrm{tol}}$-slack, then:
\[
\mathbb{E}[R_T^{\mathrm{net}}]\geq\Omega(\epsilon_{\mathrm{tol}}\,NT).
\]
\end{theorem}

\begin{proof}
Let $\beta^{(i)}(\bm{x},a)=\epsilon_{\mathrm{tol}}
\cdot\mathrm{sign}(\langle\phi(\bm{x},a),\bm{\delta}_0\rangle)$
for a fixed unit vector $\bm{\delta}_0$.
Any $\epsilon_{\mathrm{tol}}$-trust function accepts all this data.
The $\gamma$-decoupling argument (Theorem~\ref{thm:cec} with
$\gamma=0$, $\bm{\delta}=\epsilon_{\mathrm{tol}}\bm{\delta}_0$)
gives instantaneous latent regret $\Omega(\epsilon_{\mathrm{tol}})$
per round per agent, summing to $\Omega(\epsilon_{\mathrm{tol}}NT)$.
\end{proof}

\subsection{The Gossip-Error Term (iii): Open Question}

\begin{conjecture}[Gossip term refinement]\label{conj:gossip}
Term~(iii) of Theorem~\ref{thm:dcesa},
$\widetilde{O}(d\sqrt{NT}/\gamma(W))$, is not tight.  A sharper
analysis should improve the $1/\gamma(W)$ to $1/\sqrt{\gamma(W)}$,
yielding $\widetilde{O}(d\sqrt{NT/\gamma(W)})$.  (The $\sqrt N$ factor,
by contrast, is genuine, from Lemma~\ref{lem:gossip}.)
\end{conjecture}

\subsection{Summary}

\begin{center}
\begin{tabular}{lccc}
\hline
Term & Upper bound & Lower bound & Tight?\\
\hline
(i) Epistemic & $|E|\sqrt{Td/\bar{p}}$ &
$|E|\sqrt{Td/\bar{p}}$ soft-label (Thm.~\ref{thm:lower_epistemic}) &
soft: cond.; hard: open\\
(ii) Cooperative bandit & $d\sqrt{NT}$ &
$d\sqrt{NT}$ (Thm.~\ref{thm:lower_coop}) & $\checkmark$\\
(iii) Gossip error & $d\sqrt{NT}/\gamma(W)$ & open &
open (Conj.~\ref{conj:gossip})\\
(iv) Tolerance & $\epsilon_{\mathrm{tol}}NT$ &
$\epsilon_{\mathrm{tol}}NT$ (Thm.~\ref{thm:lower_tol}) & $\checkmark$\\
\hline
\end{tabular}
\end{center}

%%============================================================
\section{Remaining Caveats}\label{sec:caveats}
%%============================================================

\textbf{(C1) From population contraction to trajectory convergence
(main open problem).}
Lemma~\ref{lem:fp} gives existence; Lemma~\ref{lem:contraction} gives
\emph{uniqueness and geometric population convergence} under the
explicit contraction condition $\kappa<1$~\eqref{eq:kappa} (a margin /
low-noise / small-bias condition).  What is \emph{not} proved is that
the stochastic recursive-least-squares trajectory
$\bm{\theta}_\tau^{\mathrm{cen}}$ inherits this convergence with the
tracking partial sums $E_t=O(\sqrt t)$ of
Assumption~\ref{ass:pol_conv}; this is a stochastic-approximation
statement.  Consequently Theorem~\ref{thm:cec} is stated
\emph{conditionally} on Assumption~\ref{ass:pol_conv}
(Remark~\ref{rem:pol_conv}).  Deriving Assumption~\ref{ass:pol_conv}
from $\kappa<1$ --- i.e.\ the \emph{unconditional} supercritical
collapse theorem --- is the principal open problem.  Note the
subcritical Theorem~\ref{thm:cec_sub} is \emph{unconditional}
(it does not use Assumption~\ref{ass:pol_conv}).

\textbf{(C6) Epistemic lower bound: hard-label case open.}
Theorem~\ref{thm:lower_epistemic} establishes
$\Omega(|E|\sqrt{Td/\bar p})$ only under the soft-certification
Assumption~\ref{ass:softcert}.  For the noiseless deterministic axiom
of Section~\ref{sec:dcesa_algo}, only the weaker waiting-time bound
$\Omega(|E|\Delta_{\mathrm{gap}}/\bar p)$ is proved
(Remark~\ref{rem:epi_downgrade}); the matching $\sqrt{Td/\bar p}$
lower bound for hard labels is open.  Accordingly the earlier
``$\checkmark$ tight'' claim for term~(i) has been withdrawn.

\textbf{(C2) Doubly-stochastic $W$.}
For weighted or directed gossip, all arguments go through with
$\bar{\bm{u}}$ replaced by $\sum_i\pi_i\bm{u}_i$ for the Perron
eigenvector $\pi$.

\textbf{(C3) Intermediate regime.}
The regime $c_0/T^2<\gamma(W)<c/T$ is not covered; we conjecture
continuous interpolation from individual bias $\bm{\delta}_i$ to
network bias $\bm{\delta}$ as connectivity increases.

\textbf{(C4) Gossip term tightness.}
Conjecture~\ref{conj:gossip} states that term~(iii)
of Theorem~\ref{thm:dcesa}, $\widetilde{O}(d\sqrt{NT}/\gamma(W))$, is
improvable to $\widetilde{O}(d\sqrt{NT/\gamma(W)})$.
Adapting the self-normalized martingale argument of
Dubey--Pentland (2020) to the trust-weighted setting
would resolve this.  The $\sqrt N$ prefactor (from the corrected
gossip bound, Remark~\ref{rem:sqrtN}) is not expected to be removable.

\textbf{(C5) Propagation breakdown.}
When $K=\log N/\gamma(W)$ exceeds the horizon $T$,
Lemma~\ref{lem:prop}'s equalization of axiom rates fails,
and term~(i) of Theorem~\ref{thm:dcesa} reverts to
$\widetilde{O}(|E|\sqrt{Td/\min_i p_i})$.
A two-regime picture for the D-CESA upper bound is left for
future work.

\begin{thebibliography}{9}
\bibitem{abbasi2011}
Y.~Abbasi-Yadkori, D.~P\'al, C.~Szepesv\'ari.
\textit{Improved algorithms for linear stochastic bandits.}
NeurIPS, 2011.

\bibitem{dubey2020}
A.~Dubey, A.~Pentland.
\textit{Cooperative multi-agent bandits with heavy tails.}
ICML, 2020.

\bibitem{martinez2019}
D.~Mart\'inez-Rubio, V.~Kanade, P.~Rebeschini.
\textit{Decentralized cooperative stochastic bandits.}
NeurIPS, 2019.

\bibitem{koloskova2020}
A.~Koloskova, N.~Loizou, S.~Boreiri, M.~Jaggi, S.U.~Stich.
\textit{A unified theory of decentralized SGD with changing topology
and local updates.}
ICML, 2020.

\bibitem{ghasemi2026}
M.~Ghasemi, M.~Crowley.
\textit{Learning When to Trust in Contextual Social Bandits.}
Preprint, 2026.
\end{thebibliography}

\end{document}